\section{アミノ酸の共通骨格・側鎖・立体配置と必須性}

\subsection{平成25年度 問4-5}
\label{q:h25-reg-basic-4-5}

\noindent\textbf{原典：} sources/H25/basic.pdf、PDF物理17頁。

\textbf{公式正答：} 1（single）。

\paragraph{問題の要旨} 標準タンパク質のアミノ酸種数、共通官能基、側鎖の疎水性、立体配置を4空欄へ対応させる。

\Needspace{4\baselineskip}\paragraph{確認した条件・記述（独自要約）}\mbox{}\par

\noindent\textbf{空欄：} ア〜エ。特殊なアミノ酸を除く。\par

\noindent\textbf{対象：} タンパク質加水分解のアミノ酸、グリシンの不斉例外。\par

\paragraph{問われる核心} 共通アミノ酸骨格と側鎖・立体配置

\paragraph{必要知識} 共通アミノ酸骨格と側鎖・立体配置。原典が特殊なものを除くと限定しているため標準20種を用いる。α炭素はアミノ基・カルボキシ基・H・Rに結合する。ロイシン等の非極性Rは疎水性で、グリシンはR=\allowbreak{}Hで不斉炭素を持たず、他の標準タンパク質アミノ酸は主にLである。20・カルボキシ・疎水・Lを満たす1。D/Lは光学回転の正負ではない。

\paragraph{考え方・正答を選ぶ理由（査読待ち初稿）} 原典が特殊なものを除くと限定しているため標準20種を用いる。α炭素はアミノ基・カルボキシ基・H・Rに結合する。ロイシン等の非極性Rは疎水性で、グリシンはR=\allowbreak{}Hで不斉炭素を持たず、他の標準タンパク質アミノ酸は主にLである。20・カルボキシ・疎水・Lを満たす1。D/Lは光学回転の正負ではない。

\Needspace{4\baselineskip}\paragraph{選択肢の検討（査読待ち初稿）}\mbox{}\par

\noindent\textbf{1：最適} 20／カルボキシ基／疎水／L体。標準20種、共通カルボキシ基、非極性の疎水側鎖、主にLという全項目が一致。\par

\noindent\textbf{2：不適切} 20／ヒドロキシ基／疎水／D体。共通の基をヒドロキシ基へ置換し、標準タンパク質の立体配置も逆。\par

\noindent\textbf{3：不適切} 30／カルボキシ基／親水／L体。標準種数とロイシン・フェニルアラニンの側鎖の分類が違う。\par

\noindent\textbf{4：不適切} 30／ヒドロキシ基／親水／L体。種数・共通基・炭化水素側鎖の分類が違う。\par

\noindent\textbf{5：不適切} 30／カルボキシ基／疎水／D体。共通基と疎水性は合うが、標準20種と主にLを満たさない。\par

\paragraph{再利用する考え方} 単位・対象・結合の階層を区別する。共通アミノ酸骨格と側鎖・立体配置の成立条件を確認して類似年度の肢順へ対応する。

\paragraph{共通理論} \hyperref[theory:biology:amino]{第\ref*{theory:biology:amino}節}

\par 参考：\cite{mit_amino}。

\paragraph{関連問題} \hyperref[q:h29-reg-basic-4-5]{H29-REG-BASIC-4-5}

\par 難易度 2（校準初稿）、優先度未評価。

\subsection{平成29年度 問4-5}
\label{q:h29-reg-basic-4-5}

\noindent\textbf{原典：} sources/H29/basic.pdf、PDF物理17頁。

\textbf{公式正答：} 1（single）。

\paragraph{問題の要旨} 標準タンパク質のアミノ酸種数、共通官能基、側鎖の疎水性、立体配置を4空欄へ対応させる。

\Needspace{4\baselineskip}\paragraph{確認した条件・記述（独自要約）}\mbox{}\par

\noindent\textbf{空欄：} ア〜エ。特殊なアミノ酸を除く。\par

\noindent\textbf{対象：} タンパク質加水分解のアミノ酸、グリシンの不斉例外。\par

\paragraph{問われる核心} 共通アミノ酸骨格と側鎖・立体配置

\paragraph{必要知識} 共通アミノ酸骨格と側鎖・立体配置。原典が特殊なものを除くと限定しているため標準20種を用いる。α炭素はアミノ基・カルボキシ基・H・Rに結合する。ロイシン等の非極性Rは疎水性で、グリシンはR=\allowbreak{}Hで不斉炭素を持たず、他の標準タンパク質アミノ酸は主にLである。20・カルボキシ・疎水・Lを満たす1。D/Lは光学回転の正負ではない。

\paragraph{考え方・正答を選ぶ理由（査読待ち初稿）} 原典が特殊なものを除くと限定しているため標準20種を用いる。α炭素はアミノ基・カルボキシ基・H・Rに結合する。ロイシン等の非極性Rは疎水性で、グリシンはR=\allowbreak{}Hで不斉炭素を持たず、他の標準タンパク質アミノ酸は主にLである。20・カルボキシ・疎水・Lを満たす1。D/Lは光学回転の正負ではない。

\Needspace{4\baselineskip}\paragraph{選択肢の検討（査読待ち初稿）}\mbox{}\par

\noindent\textbf{1：最適} 20／カルボキシ基／疎水／L体。標準20種、共通カルボキシ基、非極性の疎水側鎖、主にLという全項目が一致。\par

\noindent\textbf{2：不適切} 20／ヒドロキシ基／疎水／D体。共通の基をヒドロキシ基へ置換し、標準タンパク質の立体配置も逆。\par

\noindent\textbf{3：不適切} 30／カルボキシ基／親水／L体。標準種数とロイシン・フェニルアラニンの側鎖の分類が違う。\par

\noindent\textbf{4：不適切} 30／カルボキシ基／疎水／D体。共通基と疎水性は合うが、標準20種と主にLを満たさない。\par

\noindent\textbf{5：不適切} 30／ヒドロキシ基／親水／L体。種数・共通基・炭化水素側鎖の分類が違う。\par

\paragraph{再利用する考え方} 単位・対象・結合の階層を区別する。共通アミノ酸骨格と側鎖・立体配置の成立条件を確認して類似年度の肢順へ対応する。

\paragraph{共通理論} \hyperref[theory:biology:amino]{第\ref*{theory:biology:amino}節}

\par 参考：\cite{mit_amino}。

\paragraph{関連問題} \hyperref[q:h25-reg-basic-4-5]{H25-REG-BASIC-4-5}

\par 難易度 2（校準初稿）、優先度未評価。

\subsection{令和3年度 問4-5}
\label{q:r03-reg-basic-4-5}

\noindent\textbf{原典：} sources/R03/basic.pdf、PDF物理19頁。

\textbf{公式正答：} 1（single）。

\paragraph{問題の要旨} アミノ酸の共通基、疎水側鎖、CysとMetに含まれる元素を3空欄へ対応させる。

\Needspace{4\baselineskip}\paragraph{確認した条件・記述（独自要約）}\mbox{}\par

\noindent\textbf{空欄：} ア〜ウ。Cys2側鎖間の共有結合に言及。\par

\paragraph{問われる核心} 側鎖の官能基とジスルフィド架橋

\paragraph{必要知識} 側鎖の官能基とジスルフィド架橋。α炭素に結合する共通基はカルボキシで、Ile/Pheの非極性Rは疎水性。Cysはチオール、Metはチオエーテルの硫黄を持つので1。Cysの2側鎖がS-S共有結合を作ることと、Metも同じ架橋を作ることは別である。

\paragraph{考え方・正答を選ぶ理由（査読待ち初稿）} α炭素に結合する共通基はカルボキシで、Ile/Pheの非極性Rは疎水性。Cysはチオール、Metはチオエーテルの硫黄を持つので1。Cysの2側鎖がS-S共有結合を作ることと、Metも同じ架橋を作ることは別である。

\Needspace{4\baselineskip}\paragraph{選択肢の検討（査読待ち初稿）}\mbox{}\par

\noindent\textbf{1：最適} カルボキシ基／疎水／S。共通基、Ile/Pheの非極性R、Cys/Metの硫黄が全て一致。\par

\noindent\textbf{2：不適切} ヒドロキシ基／疎水／S。RによってはOHを持つが、全アミノ酸の共通基はカルボキシである。\par

\noindent\textbf{3：不適切} カルボキシ基／親水／S。Ile/PheのRは炭化水素で疎水性である。\par

\noindent\textbf{4：不適切} カルボキシ基／親水／N。疎水性と側鎖の硫黄を取り違えている。アミノ基のNは側鎖共通元素という意味ではない。\par

\noindent\textbf{5：不適切} ヒドロキシ基／親水／N。共通基、側鎖極性、Cys/Metの特徴的元素の3項目が違う。\par

\paragraph{再利用する考え方} 単位・対象・結合の階層を区別する。側鎖の官能基とジスルフィド架橋の成立条件を確認して類似年度の肢順へ対応する。

\paragraph{共通理論} \hyperref[theory:biology:amino]{第\ref*{theory:biology:amino}節}, \hyperref[theory:biology:protein]{第\ref*{theory:biology:protein}節}

\par 参考：\cite{mit_amino,mit_protein}。

\par 難易度 2（校準初稿）、優先度未評価。

\subsection{令和7年度 問4-5}
\label{q:r07-reg-basic-4-5}

\noindent\textbf{原典：} sources/R07/basic.pdf、PDF物理14頁。

\textbf{公式正答：} 4（single）。

\paragraph{問題の要旨} 通常の必須アミノ酸の候補から、体内で合成できるものを選ぶ。

\Needspace{4\baselineskip}\paragraph{確認した条件・記述（独自要約）}\mbox{}\par

\noindent\textbf{方向：} 必須アミノ酸として最も不適切。\par

\noindent\textbf{対象：} 原典は哺乳類。解説は通常のヒトの分類と条件付き必須の差を明示。\par

\paragraph{問われる核心} アミノ酸の生合成能力と栄養上の必要条件

\paragraph{必要知識} アミノ酸の生合成能力と栄養上の必要条件。一般的なヒトの栄養分類ではIle・Val・Leu・Metは必須、Proは通常非必須なので4。必須は分子の重要性でなく、必要量を体内合成で満たせず外部供給を要するという分類である。原典は哺乳類を広く述べるが種・発育・病態の差を無条件に同一としない。Proの条件付き必須と通常の非必須を区別する。

\paragraph{考え方・正答を選ぶ理由（査読待ち初稿）} 一般的なヒトの栄養分類ではIle・Val・Leu・Metは必須、Proは通常非必須なので4。必須は分子の重要性でなく、必要量を体内合成で満たせず外部供給を要するという分類である。原典は哺乳類を広く述べるが種・発育・病態の差を無条件に同一としない。Proの条件付き必須と通常の非必須を区別する。

\Needspace{4\baselineskip}\paragraph{選択肢の検討（査読待ち初稿）}\mbox{}\par

\noindent\textbf{1：必須} イソロイシン。分岐鎖の必須アミノ酸で、通常のヒトの必須9種に含まれる。\par

\noindent\textbf{2：必須} バリン。分岐鎖の必須アミノ酸である。\par

\noindent\textbf{3：必須} ロイシン。分岐鎖の必須アミノ酸である。\par

\noindent\textbf{4：通常非必須} プロリン。通常はグルタミン酸等から生合成できる非必須に分類され、この候補が不適切。ただし病態・成長等で条件付き必須となる場合を区別。\par

\noindent\textbf{5：必須} メチオニン。硫黄を含む必須アミノ酸である。CysとMetを同じ栄養分類とする必要はない。\par

\paragraph{再利用する考え方} 単位・対象・結合の階層を区別する。アミノ酸の生合成能力と栄養上の必要条件の成立条件を確認して類似年度の肢順へ対応する。

\paragraph{共通理論} \hyperref[theory:biology:nutrition]{第\ref*{theory:biology:nutrition}節}

\par 参考：\cite{nih_amino}。

\par 難易度 2（校準初稿）、優先度未評価。

\section{機能性電子セラミックス}

\subsection{平成25年度 問4-3}
\label{q:h25-reg-basic-4-3}

\noindent\textbf{原典：} sources/H25/basic.pdf、PDF物理16頁。

\textbf{公式正答：} 4（single）。

\paragraph{問題の要旨} 電子セラミックス4種の利用特性・材料名の組合せを選ぶ。

\Needspace{4\baselineskip}\paragraph{確認した条件・記述（独自要約）}\mbox{}\par

\noindent\textbf{A：} BaTiO3の大きな何をコンデンサーへ利用するか。\par

\noindent\textbf{B：} サーミスタの温度で大きく変わる量。\par

\noindent\textbf{C：} ひずみにより電圧を生む材料。\par

\noindent\textbf{D：} バリスタの印加電圧で大きく変わる量。\par

\noindent\textbf{並び：} A／B／C／D\par

\paragraph{問われる核心} 電子セラミックスの入力・出力と物性

\paragraph{必要知識} 誘電率と導電率の違い、温度抵抗、圧電効果、非線形抵抗。

\paragraph{考え方・正答を選ぶ理由（査読待ち初稿）} コンデンサーは\(C=\varepsilon_0\varepsilon_r A/d\)で高い比誘電率を大容量へ利用する。サーミスタは抵抗の温度依存、圧電体は機械ひずみと電気応答の結合、バリスタは非線形な電圧–電流特性を用いる。比誘電率／電気抵抗／圧電体／電気抵抗の4である。

\Needspace{4\baselineskip}\paragraph{選択肢の検討（査読待ち初稿）}\mbox{}\par

\noindent\textbf{1：条件に不一致} 導電率／電気抵抗／圧電体／体積。BとCは合うが、コンデンサーは高比誘電率を利用し、バリスタは体積変化を機能としていない。\par

\noindent\textbf{2：条件に不一致} 導電率／熱膨張係数／放電体／電気抵抗。AとBは違い、機械的ひずみからの電圧発生も放電体の説明ではない。\par

\noindent\textbf{3：条件に不一致} 比誘電率／電気抵抗／放電体／体積。AとBは合うが、ひずみと電気応答は圧電体、バリスタは電気抵抗の変化である。\par

\noindent\textbf{4：条件に一致} 比誘電率／電気抵抗／圧電体／電気抵抗。4項とも対応する。誘電・温度依存抵抗・圧電・非線形抵抗を区別できている。\par

\noindent\textbf{5：条件に不一致} 比誘電率／熱膨張係数／圧電体／体積。AとCは合うが、温度で変わる抵抗と電圧で変わる抵抗を取り違える。\par

\paragraph{再利用する考え方} 電子セラミックスの入力・出力と物性を、数える対象・分母・電荷保存・成立条件を明示して判断する。

\paragraph{共通理論} \hyperref[theory:basicchem:ceramics]{第\ref*{theory:basicchem:ceramics}節}

\par 参考：\cite{murata_batio3,murata_thermistor,murata_piezo,tdk_varistor}。

\par 難易度 2（校準初稿）、優先度未評価。

\section{濃度・組成の分率}

\subsection{平成26年度 問4-2}
\label{q:h26-reg-basic-4-2}

\noindent\textbf{原典：} sources/H26/basic.pdf、PDF物理13頁。

\textbf{公式正答：} 2（single）。

\par\textbf{検討未確定：} ⑤はppmの基準と溶液密度が未指定。1 mg/L≈質量1 ppmという希薄水溶液近似の範囲はAstra確認。独立の選択2と公式2は一致する。

\paragraph{問題の要旨} ショ糖水溶液の濃度と調製方法の説明から、不適切なものを選ぶ。

\Needspace{4\baselineskip}\paragraph{確認した条件・記述（独自要約）}\mbox{}\par

\noindent\textbf{モル質量：} ショ糖342 g/mol、水18 g/mol。\par

\noindent\textbf{方向：} 不適切な説明を一つ。\par

\noindent\textbf{濃度：} 質量モル濃度、モル濃度、質量\%、モル分率、ppm。\par

\paragraph{問われる核心} 溶媒・溶液と濃度の分母

\paragraph{必要知識} 物質量、最終体積、質量モル濃度、質量分率、モル分率、密度。

\paragraph{考え方・正答を選ぶ理由（査読待ち初稿）} 定義を先に書くと、質量モル濃度は溶質mol/溶媒kg、モル濃度は溶質mol/最終溶液Lである。2だけは最終溶液体積を保証しない。1・3・4は定義式で確かめられる。5は試験の希薄水溶液・密度約1 kg/Lという近似なら妥当であり、公式の2と独立判断は一致する。ppmをmg/Lの無条件の別名として覚えない。

\Needspace{4\baselineskip}\paragraph{選択肢の検討（査読待ち初稿）}\mbox{}\par

\noindent\textbf{1：適切} 1 mol/kgはショ糖342 gを水1000 gへ溶かす。。1 molを溶媒1 kgで割る定義で、溶液全体の質量ではない。\par

\noindent\textbf{2：不適切} 1 mol/Lはショ糖342 gを水1 Lへ溶かす。。分母は最終溶液体積である。水1 Lへ加えた後の溶液が1 Lとは限らず、この調製法で濃度を保証できない。\par

\noindent\textbf{3：適切} 質量10\%はショ糖100 gを水900 gへ溶かす。。100/(100+900)=\allowbreak{}0.10である。\par

\noindent\textbf{4：適切} ショ糖のモル分率0.01は34.2 gを水178.2 gへ溶かす。。ショ糖0.1 mol、水9.9 molなので0.1/(0.1+9.9)=\allowbreak{}0.01である。\par

\noindent\textbf{5：希薄水溶液近似で適切・厳密基準は要確認} 1 ppmはショ糖1 mgを含む水溶液の体積が1 Lである。。1 mg/Lという質量濃度は確認できる。密度約1 kg/Lの希薄水溶液では質量ppmと近似一致するが、ppmの基準と密度を明記しない原典を厳密な同一量とはしない。\par

\paragraph{再利用する考え方} 溶媒・溶液と濃度の分母を、数える対象・分母・電荷保存・成立条件を明示して判断する。

\paragraph{共通理論} \hyperref[theory:basicchem:concentration]{第\ref*{theory:basicchem:concentration}節}, \hyperref[theory:basicchem:fraction]{第\ref*{theory:basicchem:fraction}節}

\par 参考：\cite{iupac_amount_concentration,iupac_molality,iupac_amount_fraction}。

\par 難易度 2（校準初稿）、優先度未評価。

\subsection{令和7年度 問4-4}
\label{q:r07-reg-basic-4-4}

\noindent\textbf{原典：} sources/R07/basic.pdf、PDF物理14頁。

\textbf{公式正答：} 5（single）。

\paragraph{問題の要旨} Sn–Ag–Cu合金の質量分率をモル分率へ換算する。

\Needspace{4\baselineskip}\paragraph{確認した条件・記述（独自要約）}\mbox{}\par

\noindent\textbf{質量分率：} Sn96.5\%、Ag3.00\%、Cu0.500\%。\par

\noindent\textbf{原子量：} Sn118.7、Ag107.8、Cu63.5。\par

\noindent\textbf{候補の並び：} Sn／Ag／Cuのモル分率\%。\par

\paragraph{問われる核心} 質量分率から物質量分率への正規化

\paragraph{必要知識} n=\allowbreak{}w/M、原子量、分率の分母、比の不変性。

\paragraph{考え方・正答を選ぶ理由（査読待ち初稿）} 合金100 gを取ると\(n_{Sn}=96.5/118.7\simeq0.8138\)、\(n_{Ag}=3.00/107.8\simeq0.02783\)、\(n_{Cu}=0.500/63.5\simeq0.007874\) molで、合計約0.8495 molである。各値を合計で割り100倍すると約95.8、3.28、0.928\%、5となる。基準質量を1 gに変えても分子・分母の係数が消えて同じ割合になる。

\Needspace{4\baselineskip}\paragraph{選択肢の検討（査読待ち初稿）}\mbox{}\par

\noindent\textbf{1：条件に不一致} 97.0／2.74／0.269。各質量を原子量で割って正規化した組と違い、特にCuが小さすぎる。\par

\noindent\textbf{2：条件に不一致} 97.0／2.34／0.662。正規化したAg約3.28\%、Cu約0.928\%と合わない。\par

\noindent\textbf{3：条件に不一致} 96.5／3.00／0.500。質量分率をそのまま書いており原子量の差が反映されない。\par

\noindent\textbf{4：条件に不一致} 95.8／3.82／0.375。Snは合うがAgとCuのモル比が与えられた質量/原子量と合わない。\par

\noindent\textbf{5：条件に一致} 95.8／3.28／0.928。各成分の物質量を合計で割った95.8\%、3.28\%、0.928\%に一致する。\par

\paragraph{再利用する考え方} 質量分率から物質量分率への正規化を、数える対象・分母・電荷保存・成立条件を明示して判断する。

\paragraph{共通理論} \hyperref[theory:basicchem:fraction]{第\ref*{theory:basicchem:fraction}節}

\par 参考：\cite{iupac_amount_fraction}。

\par 難易度 2（校準初稿）、優先度未評価。

\section{ハロゲンの物性・酸化力・酸性度}

\subsection{令和元年度 問4-1}
\label{q:r01-reg-basic-4-1}

\noindent\textbf{原典：} sources/R01/basic.pdf、PDF物理16頁。

\textbf{公式正答：} 4（single）。

\paragraph{問題の要旨} ハロゲン原子・分子・ハロゲン化水素の4記述から、正しい2項の組を選ぶ。

\Needspace{4\baselineskip}\paragraph{確認した条件・記述（独自要約）}\mbox{}\par

\noindent\textbf{ア：} ア：水溶液中の酸の強さはHF>HCl>HBr>HI。\par

\noindent\textbf{イ：} イ：原子の電気陰性度はF>Cl>Br>I。\par

\noindent\textbf{ウ：} ウ：沸点はHF>HCl>HBr>HI。\par

\noindent\textbf{エ：} エ：ハロゲン分子の酸化力はF2>Cl2>Br2>I2。\par

\paragraph{問われる核心} ハロゲン族の異なる性質を別々に比較する

\paragraph{必要知識} 電気陰性度、酸化力、酸解離と溶媒、水素結合と分散力。

\paragraph{考え方・正答を選ぶ理由（査読待ち初稿）} 電気陰性度とハロゲン分子の酸化力はF、Cl、Br、Iの順に小さくなる。しかし水中でHFは弱酸、HCl・HBr・HIは強酸であり、Fの高電気陰性度をHFが最強の酸という結論へ転用できない。強酸3種は水による水平化にも注意する。通常沸点はHF>HI>HBr>HClで、HFの水素結合と他3種の分散力の増大を分ける。正しいのはイ・エ、4である。

\Needspace{4\baselineskip}\paragraph{選択肢の検討（査読待ち初稿）}\mbox{}\par

\noindent\textbf{1：条件に不一致} ア・イ。含まれるアが誤り。酸性度・H\textsuperscript{+}濃度はHFを最大にできず、沸点ではHClが最小である。\par

\noindent\textbf{2：条件に不一致} ア・ウ。含まれるア・ウが誤り。酸性度・H\textsuperscript{+}濃度はHFを最大にできず、沸点ではHClが最小である。\par

\noindent\textbf{3：条件に不一致} イ・ウ。含まれるウが誤り。酸性度・H\textsuperscript{+}濃度はHFを最大にできず、沸点ではHClが最小である。\par

\noindent\textbf{4：条件に一致} イ・エ。両記述が正しい。\par

\noindent\textbf{5：条件に不一致} ウ・エ。含まれるウが誤り。酸性度・H\textsuperscript{+}濃度はHFを最大にできず、沸点ではHClが最小である。\par

\paragraph{再利用する考え方} ハロゲン族の異なる性質を別々に比較するを、数える対象・分母・電荷保存・成立条件を明示して判断する。

\paragraph{共通理論} \hyperref[theory:basicchem:halogen]{第\ref*{theory:basicchem:halogen}節}

\par 参考：\cite{purdue_halogens,rsc_halogen,nist_halogen_properties,msu_hydrogen_bonding,purdue_acid_strength,yale_hydrohalic}。

\paragraph{関連問題} \hyperref[q:r06-reg-basic-4-1]{R06-REG-BASIC-4-1}

\par 難易度 2（校準初稿）、優先度未評価。

\subsection{令和6年度 問4-1}
\label{q:r06-reg-basic-4-1}

\noindent\textbf{原典：} sources/R06/basic.pdf、PDF物理15頁。

\textbf{公式正答：} 1（single）。

\paragraph{問題の要旨} ハロゲン原子・分子・ハロゲン化水素の4記述から、正しい2項の組を選ぶ。

\Needspace{4\baselineskip}\paragraph{確認した条件・記述（独自要約）}\mbox{}\par

\noindent\textbf{ア：} ア：原子の電気陰性度はF>Cl>Br>I。\par

\noindent\textbf{イ：} イ：ハロゲン分子の酸化力はF2>Cl2>Br2>I2。\par

\noindent\textbf{ウ：} ウ：同濃度水溶液のH\textsuperscript{+}濃度はHF>HCl>HBr>HI。\par

\noindent\textbf{エ：} エ：沸点はHF>HCl>HBr>HI。\par

\paragraph{問われる核心} ハロゲン族の異なる性質を別々に比較する

\paragraph{必要知識} 電気陰性度、酸化力、酸解離と溶媒、水素結合と分散力。

\paragraph{考え方・正答を選ぶ理由（査読待ち初稿）} 電気陰性度とハロゲン分子の酸化力はF、Cl、Br、Iの順に小さくなる。しかし水中でHFは弱酸、HCl・HBr・HIは強酸であり、Fの高電気陰性度をHFが最強の酸という結論へ転用できない。強酸3種は水による水平化にも注意する。通常沸点はHF>HI>HBr>HClで、HFの水素結合と他3種の分散力の増大を分ける。正しいのはア・イ、1である。

\Needspace{4\baselineskip}\paragraph{選択肢の検討（査読待ち初稿）}\mbox{}\par

\noindent\textbf{1：条件に一致} ア・イ。両記述が正しい。\par

\noindent\textbf{2：条件に不一致} ア・ウ。含まれるウが誤り。酸性度・H\textsuperscript{+}濃度はHFを最大にできず、沸点ではHClが最小である。\par

\noindent\textbf{3：条件に不一致} イ・ウ。含まれるウが誤り。酸性度・H\textsuperscript{+}濃度はHFを最大にできず、沸点ではHClが最小である。\par

\noindent\textbf{4：条件に不一致} イ・エ。含まれるエが誤り。酸性度・H\textsuperscript{+}濃度はHFを最大にできず、沸点ではHClが最小である。\par

\noindent\textbf{5：条件に不一致} ウ・エ。含まれるウ・エが誤り。酸性度・H\textsuperscript{+}濃度はHFを最大にできず、沸点ではHClが最小である。\par

\paragraph{再利用する考え方} ハロゲン族の異なる性質を別々に比較するを、数える対象・分母・電荷保存・成立条件を明示して判断する。

\paragraph{共通理論} \hyperref[theory:basicchem:halogen]{第\ref*{theory:basicchem:halogen}節}

\par 参考：\cite{purdue_halogens,rsc_halogen,nist_halogen_properties,msu_hydrogen_bonding,purdue_acid_strength,yale_hydrohalic}。

\paragraph{関連問題} \hyperref[q:r01-reg-basic-4-1]{R01-REG-BASIC-4-1}

\par 難易度 2（校準初稿）、優先度未評価。

\section{酸化数・電子配分}

\subsection{平成27年度 問4-2}
\label{q:h27-reg-basic-4-2}

\noindent\textbf{原典：} sources/H27/basic.pdf、PDF物理12頁。

\textbf{公式正答：} 3（single）。

\paragraph{問題の要旨} 原典で下線を付けた原子の酸化数が最大の物質を選ぶ。

\Needspace{4\baselineskip}\paragraph{確認した条件・記述（独自要約）}\mbox{}\par

\noindent\textbf{対象：} 各肢に元素名を添え、下線対象を明記。\par

\noindent\textbf{規則：} この組ではH=+1、O=−2、単体=0。全原子の酸化数の総和は化学種の電荷。\par

\paragraph{問われる核心} 全電荷から原子1個の酸化数を決める

\paragraph{必要知識} 形式的電子配分、HとOの通常値、単体0、多原子イオンの総電荷。

\paragraph{考え方・正答を選ぶ理由（査読待ち初稿）} 化学種全体の電荷を右辺にして未知原子の酸化数xを解く。例えばMnO4\textsuperscript{-}はx+4(−2)=−1よりx=+7、Cr2O7\textsuperscript{2-}は2x+7(−2)=−2よりx=+6である。NH3ならx+3(+1)=\allowbreak{}0より−3となる。この問題の並びは+5, +6, +7, +4, +6なので3である。二つのCrの総和+12とCr1原子の+6を区別する。

\Needspace{4\baselineskip}\paragraph{選択肢の検討（査読待ち初稿）}\mbox{}\par

\noindent\textbf{1：条件に不一致} HNO3のN。HNO3のNは酸化数+5。選択対象の最大値ではない。\par

\noindent\textbf{2：条件に不一致} Cr2O7\textsuperscript{2-}のCr。Cr2O7\textsuperscript{2-}のCrは酸化数+6。選択対象の最大値ではない。\par

\noindent\textbf{3：条件に一致} MnO4\textsuperscript{-}のMn。MnO4\textsuperscript{-}のMnは酸化数+7。比較した最大値である。\par

\noindent\textbf{4：条件に不一致} NO2のN。NO2のNは酸化数+4。選択対象の最大値ではない。\par

\noindent\textbf{5：条件に不一致} H2SO4のS。H2SO4のSは酸化数+6。選択対象の最大値ではない。\par

\paragraph{再利用する考え方} 全電荷から原子1個の酸化数を決めるを、数える対象・分母・電荷保存・成立条件を明示して判断する。

\paragraph{共通理論} \hyperref[theory:basicchem:oxidation]{第\ref*{theory:basicchem:oxidation}節}

\par 参考：\cite{iupac_oxidation_state}。

\paragraph{関連問題} \hyperref[q:r07-reg-basic-4-1]{R07-REG-BASIC-4-1}

\par 難易度 2（校準初稿）、優先度未評価。

\subsection{令和4年度 問4-2}
\label{q:r04-reg-basic-4-2}

\noindent\textbf{原典：} sources/R04/basic.pdf、PDF物理14頁。

\textbf{公式正答：} 4（single）。

\paragraph{問題の要旨} 原典で下線を付けた原子の酸化数が最小の物質を選ぶ。

\Needspace{4\baselineskip}\paragraph{確認した条件・記述（独自要約）}\mbox{}\par

\noindent\textbf{対象：} 各肢に元素名を添え、下線対象を明記。\par

\noindent\textbf{規則：} この組ではH=+1、O=−2、単体=0。全原子の酸化数の総和は化学種の電荷。\par

\paragraph{問われる核心} 全電荷から原子1個の酸化数を決める

\paragraph{必要知識} 形式的電子配分、HとOの通常値、単体0、多原子イオンの総電荷。

\paragraph{考え方・正答を選ぶ理由（査読待ち初稿）} 化学種全体の電荷を右辺にして未知原子の酸化数xを解く。例えばMnO4\textsuperscript{-}はx+4(−2)=−1よりx=+7、Cr2O7\textsuperscript{2-}は2x+7(−2)=−2よりx=+6である。NH3ならx+3(+1)=\allowbreak{}0より−3となる。この問題の並びは-2, 0, +7, -3, +5なので4である。二つのCrの総和+12とCr1原子の+6を区別する。

\Needspace{4\baselineskip}\paragraph{選択肢の検討（査読待ち初稿）}\mbox{}\par

\noindent\textbf{1：条件に不一致} H2SのS。H2SのSは酸化数-2。選択対象の最小値ではない。\par

\noindent\textbf{2：条件に不一致} 単体Mn。単体Mnは酸化数0。選択対象の最小値ではない。\par

\noindent\textbf{3：条件に不一致} MnO4\textsuperscript{-}のMn。MnO4\textsuperscript{-}のMnは酸化数+7。選択対象の最小値ではない。\par

\noindent\textbf{4：条件に一致} NH3のN。NH3のNは酸化数-3。比較した最小値である。\par

\noindent\textbf{5：条件に不一致} HNO3のN。HNO3のNは酸化数+5。選択対象の最小値ではない。\par

\paragraph{再利用する考え方} 全電荷から原子1個の酸化数を決めるを、数える対象・分母・電荷保存・成立条件を明示して判断する。

\paragraph{共通理論} \hyperref[theory:basicchem:oxidation]{第\ref*{theory:basicchem:oxidation}節}

\par 参考：\cite{iupac_oxidation_state}。

\par 難易度 2（校準初稿）、優先度未評価。

\subsection{令和7年度 問4-1}
\label{q:r07-reg-basic-4-1}

\noindent\textbf{原典：} sources/R07/basic.pdf、PDF物理13頁。

\textbf{公式正答：} 2（single）。

\paragraph{問題の要旨} 原典で下線を付けた原子の酸化数が最大の物質を選ぶ。

\Needspace{4\baselineskip}\paragraph{確認した条件・記述（独自要約）}\mbox{}\par

\noindent\textbf{対象：} 各肢に元素名を添え、下線対象を明記。\par

\noindent\textbf{規則：} この組ではH=+1、O=−2、単体=0。全原子の酸化数の総和は化学種の電荷。\par

\paragraph{問われる核心} 全電荷から原子1個の酸化数を決める

\paragraph{必要知識} 形式的電子配分、HとOの通常値、単体0、多原子イオンの総電荷。

\paragraph{考え方・正答を選ぶ理由（査読待ち初稿）} 化学種全体の電荷を右辺にして未知原子の酸化数xを解く。例えばMnO4\textsuperscript{-}はx+4(−2)=−1よりx=+7、Cr2O7\textsuperscript{2-}は2x+7(−2)=−2よりx=+6である。NH3ならx+3(+1)=\allowbreak{}0より−3となる。この問題の並びは+6, +7, +6, +5, +4なので2である。二つのCrの総和+12とCr1原子の+6を区別する。

\Needspace{4\baselineskip}\paragraph{選択肢の検討（査読待ち初稿）}\mbox{}\par

\noindent\textbf{1：条件に不一致} H2SO4のS。H2SO4のSは酸化数+6。選択対象の最大値ではない。\par

\noindent\textbf{2：条件に一致} MnO4\textsuperscript{-}のMn。MnO4\textsuperscript{-}のMnは酸化数+7。比較した最大値である。\par

\noindent\textbf{3：条件に不一致} Cr2O7\textsuperscript{2-}のCr。Cr2O7\textsuperscript{2-}のCrは酸化数+6。選択対象の最大値ではない。\par

\noindent\textbf{4：条件に不一致} HNO3のN。HNO3のNは酸化数+5。選択対象の最大値ではない。\par

\noindent\textbf{5：条件に不一致} NO2のN。NO2のNは酸化数+4。選択対象の最大値ではない。\par

\paragraph{再利用する考え方} 全電荷から原子1個の酸化数を決めるを、数える対象・分母・電荷保存・成立条件を明示して判断する。

\paragraph{共通理論} \hyperref[theory:basicchem:oxidation]{第\ref*{theory:basicchem:oxidation}節}

\par 参考：\cite{iupac_oxidation_state}。

\paragraph{関連問題} \hyperref[q:h27-reg-basic-4-2]{H27-REG-BASIC-4-2}

\par 難易度 2（校準初稿）、優先度未評価。

\section{結合エネルギー・反応熱}

\subsection{平成23年度 問4-2}
\label{q:h23-reg-basic-4-2}

\noindent\textbf{原典：} sources/H23/basic.pdf、PDF物理12頁。

\textbf{公式正答：} 1（single）。

\paragraph{問題の要旨} 水素と塩素からHClを1 mol生成する反応の熱を、結合エネルギーから求める。

\Needspace{4\baselineskip}\paragraph{確認した条件・記述（独自要約）}\mbox{}\par

\noindent\textbf{切断エネルギー：} H–H:436、Cl–Cl:243、H–Cl:432 kJ/mol\par

\noindent\textbf{基準：} HCl生成量1 mol。気体分子の結合エネルギーによる評価。\par

\noindent\textbf{符号：} 発生する熱の量を問う。候補は正の値。\par

\paragraph{問われる核心} 結合切断・生成の差と反応量基準

\paragraph{必要知識} ヘスの法則、気相結合エンタルピー、反応係数、発熱の符号。

\paragraph{考え方・正答を選ぶ理由（査読待ち初稿）} H2+Cl2→2HClを基準に、切断へ436+243=\allowbreak{}679 kJ、生成で2×432=\allowbreak{}864 kJが放出される。\(\Delta H=(679-864)/2=-92.5\) kJ/mol-HClである。原典の「発生する熱」と正の候補に従えば放出熱の大きさ92.5≃93の1である。エンタルピー変化そのものは負と区別する。

\Needspace{4\baselineskip}\paragraph{選択肢の検討（査読待ち初稿）}\mbox{}\par

\noindent\textbf{1：条件に一致} 93 kJ/mol。HCl1 mol当たりの放出熱92.5に最も近い。\par

\noindent\textbf{2：条件に不一致} 216 kJ/mol。切断と生成の差をHCl2 molで割った92.5の大きさに合わない。\par

\noindent\textbf{3：条件に不一致} 340 kJ/mol。切断エネルギー679の約半分だけを数え、生成側を差し引いていない。\par

\noindent\textbf{4：条件に不一致} 432 kJ/mol。HClの結合を切るエネルギーであり、反応の正味熱ではない。\par

\noindent\textbf{5：条件に不一致} 679 kJ/mol。水素と塩素の切断だけの679を使い、HCl生成による放出と1 mol基準を無視する。\par

\paragraph{再利用する考え方} 結合切断・生成の差と反応量基準を、数える対象・分母・電荷保存・成立条件を明示して判断する。

\paragraph{共通理論} \hyperref[theory:basicchem:enthalpy]{第\ref*{theory:basicchem:enthalpy}節}

\par 参考：\cite{mit_bond_enthalpy}。

\paragraph{関連問題} \hyperref[q:h26-reg-basic-4-1]{H26-REG-BASIC-4-1}

\par 難易度 2（校準初稿）、優先度未評価。

\subsection{平成26年度 問4-1}
\label{q:h26-reg-basic-4-1}

\noindent\textbf{原典：} sources/H26/basic.pdf、PDF物理13頁。

\textbf{公式正答：} 1（single）。

\paragraph{問題の要旨} 水素と塩素からHClを1 mol生成する反応の熱を、結合エネルギーから求める。

\Needspace{4\baselineskip}\paragraph{確認した条件・記述（独自要約）}\mbox{}\par

\noindent\textbf{切断エネルギー：} H–H:436、Cl–Cl:243、H–Cl:432 kJ/mol\par

\noindent\textbf{基準：} HCl生成量1 mol。気体分子の結合エネルギーによる評価。\par

\noindent\textbf{符号：} 発熱を負とする指定あり。\par

\paragraph{問われる核心} 結合切断・生成の差と反応量基準

\paragraph{必要知識} ヘスの法則、気相結合エンタルピー、反応係数、発熱の符号。

\paragraph{考え方・正答を選ぶ理由（査読待ち初稿）} H2+Cl2→2HClを基準に、切断へ436+243=\allowbreak{}679 kJ、生成で2×432=\allowbreak{}864 kJが放出される。\(\Delta H=(679-864)/2=-92.5\) kJ/mol-HClである。原典は発熱を負と指定するので−93の1である。

\Needspace{4\baselineskip}\paragraph{選択肢の検討（査読待ち初稿）}\mbox{}\par

\noindent\textbf{1：条件に一致} -93 kJ/mol。HCl1 mol当たりのエンタルピー変化−92.5に最も近い。\par

\noindent\textbf{2：条件に不一致} -216 kJ/mol。切断と生成の差をHCl2 molで割った92.5の大きさに合わない。\par

\noindent\textbf{3：条件に不一致} -340 kJ/mol。切断エネルギー679の約半分だけを数え、生成側を差し引いていない。\par

\noindent\textbf{4：条件に不一致} -432 kJ/mol。HClの結合を切るエネルギーであり、反応の正味熱ではない。\par

\noindent\textbf{5：条件に不一致} -679 kJ/mol。水素と塩素の切断だけの679を使い、HCl生成による放出と1 mol基準を無視する。\par

\paragraph{再利用する考え方} 結合切断・生成の差と反応量基準を、数える対象・分母・電荷保存・成立条件を明示して判断する。

\paragraph{共通理論} \hyperref[theory:basicchem:enthalpy]{第\ref*{theory:basicchem:enthalpy}節}

\par 参考：\cite{mit_bond_enthalpy}。

\paragraph{関連問題} \hyperref[q:h23-reg-basic-4-2]{H23-REG-BASIC-4-2}

\par 難易度 2（校準初稿）、優先度未評価。

\section{化学量論・反応収率}

\subsection{平成27年度 問4-1}
\label{q:h27-reg-basic-4-1}

\noindent\textbf{原典：} sources/H27/basic.pdf、PDF物理12頁。

\textbf{公式正答：} 4（single）。

\paragraph{問題の要旨} 5物質をそれぞれ同じ質量だけ完全燃焼させ、CO2発生量が最大のものを選ぶ。

\Needspace{4\baselineskip}\paragraph{確認した条件・記述（独自要約）}\mbox{}\par

\noindent\textbf{条件：} CO2と水への完全燃焼。同質量。\par

\noindent\textbf{原子量：} H=\allowbreak{}1、C=\allowbreak{}12、O=\allowbreak{}16。\par

\paragraph{問われる核心} 完全燃焼の炭素保存と単位質量当たり収率

\paragraph{必要知識} n=\allowbreak{}w/M、原子数の保存、分子式、完全燃焼。

\paragraph{考え方・正答を選ぶ理由（査読待ち初稿）} 分子中の炭素原子数をc、モル質量をM、燃焼させる質量をwとすると、炭素保存から\(n_{CO_2}=c w/M\)である。同じwなのでc/Mを比べればよい。1/16、1/32、2/30、2/28、2/46の最大はエチレン2/28で4を選ぶ。炭素数だけではなく、その質量を運ぶ水素・酸素も分母へ含める。

\Needspace{4\baselineskip}\paragraph{選択肢の検討（査読待ち初稿）}\mbox{}\par

\noindent\textbf{1：条件に不一致} メタン CH4。1 g当たりのCO2は1/16=\allowbreak{}0.06250 molである。エチレンの2/28より小さい。\par

\noindent\textbf{2：条件に不一致} メタノール CH4O。1 g当たりのCO2は1/32=\allowbreak{}0.03125 molである。エチレンの2/28より小さい。\par

\noindent\textbf{3：条件に不一致} エタン C2H6。1 g当たりのCO2は2/30=\allowbreak{}0.06667 molである。エチレンの2/28より小さい。\par

\noindent\textbf{4：条件に一致} エチレン C2H4。1 g当たりのCO2は2/28=\allowbreak{}0.07143 molである。5物質で最大。\par

\noindent\textbf{5：条件に不一致} エタノール C2H6O。1 g当たりのCO2は2/46=\allowbreak{}0.04348 molである。エチレンの2/28より小さい。\par

\paragraph{再利用する考え方} 完全燃焼の炭素保存と単位質量当たり収率を、数える対象・分母・電荷保存・成立条件を明示して判断する。

\paragraph{共通理論} \hyperref[theory:basicchem:amount]{第\ref*{theory:basicchem:amount}節}, \hyperref[theory:basicchem:stoichiometry]{第\ref*{theory:basicchem:stoichiometry}節}

\paragraph{関連問題} \hyperref[q:r02-reg-basic-4-1]{R02-REG-BASIC-4-1}

\par 難易度 2（校準初稿）、優先度未評価。

\subsection{平成30年度 問4-1}
\label{q:h30-reg-basic-4-1}

\noindent\textbf{原典：} sources/H30/basic.pdf、PDF物理15頁。

\textbf{公式正答：} 5（single）。

\par\textbf{検討未確定：} ④のCu酸化生成物が未指定。CuO想定0.503 mol、Cu2O想定0.2516 molで、どちらでも⑤0.250 molが最小。反応条件の欠落はAstraで確認。

\paragraph{問題の要旨} 5条件に対応する物質量を比較して、最小のものを選ぶ。

\Needspace{4\baselineskip}\paragraph{確認した条件・記述（独自要約）}\mbox{}\par

\noindent\textbf{原子量：} C12.0、H1.0、O16.0、Na23.0、Cl35.5、Cu63.6。\par

\noindent\textbf{気体条件：} 0℃、1.013×10\textsuperscript{5} Pa。理想気体モル体積約22.4 L/mol。\par

\paragraph{問われる核心} 異なる記述から物質量へそろえて比較する

\paragraph{必要知識} n=\allowbreak{}w/M、n=\allowbreak{}N/NA、標準状態の気体体積、質量分率、反応係数。

\paragraph{考え方・正答を選ぶ理由（査読待ち初稿）} 全てをmolへ統一する。気体体積、溶液の溶質質量、分子個数、反応係数を使った後は0.625、0.342、0.498、CuOなら0.503、0.250 molとなり5である。4は通常のCuOを想定して解けるが、原典に生成物の式はない。Cu2Oの場合も64/(4×63.6)≃0.2516 molで順位は変わらない。酸化物の未指定を黙って確定情報へ変えない。

\Needspace{4\baselineskip}\paragraph{選択肢の検討（査読待ち初稿）}\mbox{}\par

\noindent\textbf{1：条件に不一致} 標準状態のN2 14 L。14/22.4=\allowbreak{}0.625 molで、0.250 molより大きい。\par

\noindent\textbf{2：条件に不一致} 質量10\%の食塩水200 g中のNaCl。NaCl20 g、20/58.5≃0.342 molである。\par

\noindent\textbf{3：条件に不一致} 水分子3.0×10\textsuperscript{23}個に相当する水。3.0×10\textsuperscript{23}/(6.02214076×10\textsuperscript{23})≃0.498 molである。\par

\noindent\textbf{4：条件に不一致} Cu 64 gの酸化に消費するO2。CuO生成ならO2は64/(2×63.6)≃0.503 mol。Cu2O生成なら約0.252 molで、いずれも⑤の0.250 molより大きい。ただし生成酸化物は原典で未指定。\par

\noindent\textbf{5：条件に一致} CH4 4.0 gを完全燃焼して発生するCO2。CH4→CO2は1:1。4.0/16.0=\allowbreak{}0.250 molで最小。\par

\paragraph{再利用する考え方} 異なる記述から物質量へそろえて比較するを、数える対象・分母・電荷保存・成立条件を明示して判断する。

\paragraph{共通理論} \hyperref[theory:basicchem:amount]{第\ref*{theory:basicchem:amount}節}, \hyperref[theory:basicchem:stoichiometry]{第\ref*{theory:basicchem:stoichiometry}節}

\par 難易度 2（校準初稿）、優先度未評価。

\subsection{令和2年度 問4-1}
\label{q:r02-reg-basic-4-1}

\noindent\textbf{原典：} sources/R02/basic.pdf、PDF物理15頁。

\textbf{公式正答：} 2（single）。

\paragraph{問題の要旨} 5物質をそれぞれ同じ質量だけ完全燃焼させ、CO2発生量が最大のものを選ぶ。

\Needspace{4\baselineskip}\paragraph{確認した条件・記述（独自要約）}\mbox{}\par

\noindent\textbf{条件：} CO2と水への完全燃焼。同質量。\par

\noindent\textbf{原子量：} H=\allowbreak{}1、C=\allowbreak{}12、O=\allowbreak{}16。\par

\paragraph{問われる核心} 完全燃焼の炭素保存と単位質量当たり収率

\paragraph{必要知識} n=\allowbreak{}w/M、原子数の保存、分子式、完全燃焼。

\paragraph{考え方・正答を選ぶ理由（査読待ち初稿）} 分子中の炭素原子数をc、モル質量をM、燃焼させる質量をwとすると、炭素保存から\(n_{CO_2}=c w/M\)である。同じwなのでc/Mを比べればよい。1/16、1/32、2/30、2/28、2/46の最大はエチレン2/28で2を選ぶ。炭素数だけではなく、その質量を運ぶ水素・酸素も分母へ含める。

\Needspace{4\baselineskip}\paragraph{選択肢の検討（査読待ち初稿）}\mbox{}\par

\noindent\textbf{1：条件に不一致} メタン CH4。1 g当たりのCO2は1/16=\allowbreak{}0.06250 molである。エチレンの2/28より小さい。\par

\noindent\textbf{2：条件に一致} エチレン C2H4。1 g当たりのCO2は2/28=\allowbreak{}0.07143 molである。5物質で最大。\par

\noindent\textbf{3：条件に不一致} エタン C2H6。1 g当たりのCO2は2/30=\allowbreak{}0.06667 molである。エチレンの2/28より小さい。\par

\noindent\textbf{4：条件に不一致} メタノール CH4O。1 g当たりのCO2は1/32=\allowbreak{}0.03125 molである。エチレンの2/28より小さい。\par

\noindent\textbf{5：条件に不一致} エタノール C2H6O。1 g当たりのCO2は2/46=\allowbreak{}0.04348 molである。エチレンの2/28より小さい。\par

\paragraph{再利用する考え方} 完全燃焼の炭素保存と単位質量当たり収率を、数える対象・分母・電荷保存・成立条件を明示して判断する。

\paragraph{共通理論} \hyperref[theory:basicchem:amount]{第\ref*{theory:basicchem:amount}節}, \hyperref[theory:basicchem:stoichiometry]{第\ref*{theory:basicchem:stoichiometry}節}

\paragraph{関連問題} \hyperref[q:h27-reg-basic-4-1]{H27-REG-BASIC-4-1}

\par 難易度 2（校準初稿）、優先度未評価。

\subsection{令和6年度 問4-2}
\label{q:r06-reg-basic-4-2}

\noindent\textbf{原典：} sources/R06/basic.pdf、PDF物理15頁。

\textbf{公式正答：} 3（single）。

\paragraph{問題の要旨} 2価の多原子イオンを含む化合物を選ぶ。

\Needspace{4\baselineskip}\paragraph{確認した条件・記述（独自要約）}\mbox{}\par

\noindent\textbf{条件：} 電荷の絶対値が2で、複数原子からなるイオンを探す。化合物中に同じ1価イオンが2個あることとは違う。\par

\paragraph{問われる核心} イオン1個の組成・電荷と式単位の個数

\paragraph{必要知識} 単原子・多原子イオン、電気的中性、硫酸・硝酸・リン酸、括弧添字。

\paragraph{考え方・正答を選ぶ理由（査読待ち初稿）} まず単原子か多原子かを分け、次に式単位の電気的中性からイオン電荷を確認する。(NH4)2SO4ではNH4\textsuperscript{+}が2個なのでSO4は−2であり、3を選ぶ。括弧の外の添字2はイオン個数を示し、NO3そのものを2価に変えるものではない。

\Needspace{4\baselineskip}\paragraph{選択肢の検討（査読待ち初稿）}\mbox{}\par

\noindent\textbf{1：条件に不一致} 硫化銀 Ag2S。S\textsuperscript{2-}は2価だが単原子であり条件を満たさない。\par

\noindent\textbf{2：条件に不一致} 水酸化ナトリウム NaOH。OH\textsuperscript{-}は多原子だが1価である。\par

\noindent\textbf{3：条件に一致} 硫酸アンモニウム (NH4)2SO4。SO4\textsuperscript{2-}はSとOからなる多原子・2価陰イオンである。NH4\textsuperscript{+}は1価陽イオン。\par

\noindent\textbf{4：条件に不一致} 硝酸鉛(II) Pb(NO3)2。NO3\textsuperscript{-}は多原子・1価で、Pb\textsuperscript{2+}は2価だが単原子である。\par

\noindent\textbf{5：条件に不一致} リン酸カルシウム Ca3(PO4)2。PO4\textsuperscript{3-}は多原子・3価。Ca\textsuperscript{2+}は2価だが単原子である。\par

\paragraph{再利用する考え方} イオン1個の組成・電荷と式単位の個数を、数える対象・分母・電荷保存・成立条件を明示して判断する。

\paragraph{共通理論} \hyperref[theory:basicchem:ion]{第\ref*{theory:basicchem:ion}節}

\par 難易度 2（校準初稿）、優先度未評価。

\section{細胞組成の分母と代表値}

\subsection{平成26年度 問4-6}
\label{q:h26-reg-basic-4-6}

\noindent\textbf{原典：} sources/H26/basic.pdf、PDF物理15頁。

\textbf{公式正答：} 1（single）。

\par\textbf{独立解答と公式正答に相違あり：} 独立解答 1|2。初回1|2（主候補1）対公式1。1のH/O/P/CはNの代替誤りを確認した。2の全生物で元素組成比同様という断定は典型例の近似と普遍性を区別する必要があり、全肢確定としない。

\paragraph{問題の要旨} 原子数を分母とする主要元素、細胞の水の質量分率、小分子と核酸の単位を比較する。

\Needspace{4\baselineskip}\paragraph{確認した条件・記述（独自要約）}\mbox{}\par

\noindent\textbf{方向：} 最も不適切。\par

\noindent\textbf{原子数：} 1は99パーセント、2は60パーセント以上。\par

\noindent\textbf{質量：} 3は細菌細胞全体の約70パーセント。\par

\noindent\textbf{断定：} 2は全ての生物で同様と記載。H30/R07のよく似ておりとは区別。\par

\paragraph{問われる核心} 細胞組成の分母と代表値の適用範囲

\paragraph{必要知識} 細胞組成の分母と代表値の適用範囲。主要4元素のNをPへ替えた1が明らかな誤りで、公式は1。初回独立判断は1を主候補としつつ、2の「すべての生物で同様」にも疑義があり1|2を保存した。水の質量割合と水素の原子数割合を混同しない。典型的な細胞の概数を全生物の普遍的な同一比率へ拡張できず、2の確認を残す。

\paragraph{考え方・正答を選ぶ理由（査読待ち初稿）} 主要4元素のNをPへ替えた1が明らかな誤りで、公式は1。初回独立判断は1を主候補としつつ、2の「すべての生物で同様」にも疑義があり1|2を保存した。水の質量割合と水素の原子数割合を混同しない。典型的な細胞の概数を全生物の普遍的な同一比率へ拡張できず、2の確認を残す。

\Needspace{4\baselineskip}\paragraph{選択肢の検討（査読待ち初稿）}\mbox{}\par

\noindent\textbf{1：不適切} 原子数99パーセントはH・O・P・C。。主要4元素で通常挙げるのはH・O・N・Cである。リンは核酸・リン脂質に重要だが窒素の代わりの主要4元素ではない。\par

\noindent\textbf{2：範囲に疑義} 全ての生物で元素組成比は同様で、原子数の60パーセント以上がH。。原子数のHの大きな割合は水と有機物で説明できる。しかし全生物の比が同様という断定の範囲は原典で限定されず、完全な普遍則にはできない。\par

\noindent\textbf{3：適切（代表値）} 細菌細胞の質量約70パーセントは水。。通常の細菌細胞についての代表的な概数である。乾燥重量の70パーセントという意味ではない。\par

\noindent\textbf{4：適切} 主要な有機小分子は糖・アミノ酸・脂肪酸・ヌクレオチド。。重合体の部品や反応・エネルギー源となる4群である。全ての小分子を4群だけに限定する定義ではない。\par

\noindent\textbf{5：適切} ヌクレオチドは核酸の構成単位。。塩基・糖・リン酸からなる単位がホスホジエステル結合で核酸鎖を作る。\par

\paragraph{再利用する考え方} 単位・対象・結合の階層を区別する。細胞組成の分母と代表値の適用範囲の成立条件を確認して類似年度の肢順へ対応する。

\paragraph{共通理論} \hyperref[theory:biology:composition]{第\ref*{theory:biology:composition}節}

\par 参考：\cite{bio_composition}。

\paragraph{関連問題} \hyperref[q:h30-reg-basic-4-5]{H30-REG-BASIC-4-5}

\par 難易度 2（校準初稿）、優先度未評価。

\subsection{平成30年度 問4-5}
\label{q:h30-reg-basic-4-5}

\noindent\textbf{原典：} sources/H30/basic.pdf、PDF物理16頁。

\textbf{公式正答：} 5（single）。

\paragraph{問題の要旨} 細胞の元素組成・水と有機物の質量比を、原子数割合と区別して比較する。

\Needspace{4\baselineskip}\paragraph{確認した条件・記述（独自要約）}\mbox{}\par

\noindent\textbf{方向：} 最も不適切。\par

\noindent\textbf{比較：} 細菌の水、細胞の原子数、動物細胞有機成分の重量。\par

\paragraph{問われる核心} 細胞組成の個数・湿重量・有機物重量の分母

\paragraph{必要知識} 細胞組成の個数・湿重量・有機物重量の分母。水約70は全細胞の重量、H約60以上は原子数、核酸とタンパク質の比較は有機成分の重量である。通常の細胞でタンパク質より核酸の重量比が大きいとする5が不適切。代表値を全種・全組織の固定組成にはしない。

\paragraph{考え方・正答を選ぶ理由（査読待ち初稿）} 水約70は全細胞の重量、H約60以上は原子数、核酸とタンパク質の比較は有機成分の重量である。通常の細胞でタンパク質より核酸の重量比が大きいとする5が不適切。代表値を全種・全組織の固定組成にはしない。

\Needspace{4\baselineskip}\paragraph{選択肢の検討（査読待ち初稿）}\mbox{}\par

\noindent\textbf{1：適切（代表値）} 細菌細胞の総重量約70パーセントは水。。水を含む全体の代表的な質量比である。乾燥重量と混同しない。\par

\noindent\textbf{2：適切（代表値）} 原子数99パーセントはH・O・N・C。。典型的な細胞における原子個数の概数である。質量99パーセントとは違う。\par

\noindent\textbf{3：適切（近似）} 生物の元素組成はよく似ており、原子数のHは60パーセント以上。。原子数では軽いHの数が多い。代表的な組成の近似であり、全生物の厳密に同じ割合を保証する式ではない。\par

\noindent\textbf{4：適切} 主な有機小分子は糖・アミノ酸・脂肪酸・ヌクレオチド。。高分子の部品・反応成分としての主な4群である。\par

\noindent\textbf{5：不適切} 動物細胞の有機物の重量比で核酸が最大。。通常の動物細胞の有機成分ではタンパク質の重量が大きい。核酸が最大という比較が違う。\par

\paragraph{再利用する考え方} 単位・対象・結合の階層を区別する。細胞組成の個数・湿重量・有機物重量の分母の成立条件を確認して類似年度の肢順へ対応する。

\paragraph{共通理論} \hyperref[theory:biology:composition]{第\ref*{theory:biology:composition}節}

\par 参考：\cite{bio_composition}。

\paragraph{関連問題} \hyperref[q:h26-reg-basic-4-6]{H26-REG-BASIC-4-6}, \hyperref[q:r07-reg-basic-4-6]{R07-REG-BASIC-4-6}

\par 難易度 2（校準初稿）、優先度未評価。

\subsection{令和7年度 問4-6}
\label{q:r07-reg-basic-4-6}

\noindent\textbf{原典：} sources/R07/basic.pdf、PDF物理14頁。

\textbf{公式正答：} 4（single）。

\paragraph{問題の要旨} 人体の原子数割合、細菌の水の重量と巨大分子の重量比を比較する。

\Needspace{4\baselineskip}\paragraph{確認した条件・記述（独自要約）}\mbox{}\par

\noindent\textbf{方向：} 最も不適切。\par

\noindent\textbf{比較：} 1/2は人体原子数、4/5は細菌の巨大分子重量又は総重量。\par

\paragraph{問われる核心} 細胞組成の個数・湿重量・巨大分子重量の分母

\paragraph{必要知識} 細胞組成の個数・湿重量・巨大分子重量の分母。原子数の分母と重量の分母を区別する。標準的な細菌細胞の巨大分子の重量で最大なのはタンパク質で、核酸とする4が不適切。H30は動物細胞の有機物、今回は細菌の巨大分子という比較対象の差も記録する。

\paragraph{考え方・正答を選ぶ理由（査読待ち初稿）} 原子数の分母と重量の分母を区別する。標準的な細菌細胞の巨大分子の重量で最大なのはタンパク質で、核酸とする4が不適切。H30は動物細胞の有機物、今回は細菌の巨大分子という比較対象の差も記録する。

\Needspace{4\baselineskip}\paragraph{選択肢の検討（査読待ち初稿）}\mbox{}\par

\noindent\textbf{1：適切（代表値）} 人体原子数の90パーセント以上はH・O・N・C。。典型的な細胞における原子個数の概数である。質量90以上パーセントとは違う。\par

\noindent\textbf{2：適切（近似）} 生物の元素組成はよく似ており、人体原子数のHは60パーセント以上。。原子数では軽いHの数が多い。代表的な組成の近似であり、全生物の厳密に同じ割合を保証する式ではない。\par

\noindent\textbf{3：適切} 主な有機小分子は糖・アミノ酸・脂肪酸・ヌクレオチド。。高分子の部品・反応成分としての主な4群である。\par

\noindent\textbf{4：不適切} 細菌の巨大分子で核酸の重量比が最大。。通常の細菌の巨大分子ではタンパク質の重量が最大である。DNAとRNAを合算しても常に最大とはならない。\par

\noindent\textbf{5：適切（代表値）} 細菌細胞の総重量約70パーセントは水。。水を含む全体の代表的な質量比である。乾燥重量と混同しない。\par

\paragraph{再利用する考え方} 単位・対象・結合の階層を区別する。細胞組成の個数・湿重量・巨大分子重量の分母の成立条件を確認して類似年度の肢順へ対応する。

\paragraph{共通理論} \hyperref[theory:biology:composition]{第\ref*{theory:biology:composition}節}

\par 参考：\cite{bio_composition}。

\paragraph{関連問題} \hyperref[q:h30-reg-basic-4-5]{H30-REG-BASIC-4-5}

\par 難易度 2（校準初稿）、優先度未評価。

\section{植物の再分化と体細胞核移植}

\subsection{平成27年度 問4-5}
\label{q:h27-reg-basic-4-5}

\noindent\textbf{原典：} sources/H27/basic.pdf、PDF物理14頁。

\textbf{公式正答：} 2（single）。

\paragraph{問題の要旨} 植物の組織培養と哺乳類の体細胞核移植に必要な因子・細胞群・移植対象を対応させる。

\Needspace{4\baselineskip}\paragraph{確認した条件・記述（独自要約）}\mbox{}\par

\noindent\textbf{空欄：} ア〜エ。カルス培養後の再分化、体細胞核移植。\par

\paragraph{問われる核心} 植物の再分化制御と体細胞核の再プログラム化

\paragraph{必要知識} 植物の再分化制御と体細胞核の再プログラム化。植物ではホルモン条件を変えて細胞増殖と再分化を制御する。設問の細胞塊はカルスで、プロトプラストとは定義が違う。体細胞クローンは除核した未受精卵へ体細胞核を移植するので2。核DNAの同一性を、卵由来のミトコンドリアDNA・エピゲノム・表現型まで全て同じという主張に拡張しない。

\paragraph{考え方・正答を選ぶ理由（査読待ち初稿）} 植物ではホルモン条件を変えて細胞増殖と再分化を制御する。設問の細胞塊はカルスで、プロトプラストとは定義が違う。体細胞クローンは除核した未受精卵へ体細胞核を移植するので2。核DNAの同一性を、卵由来のミトコンドリアDNA・エピゲノム・表現型まで全て同じという主張に拡張しない。

\Needspace{4\baselineskip}\paragraph{選択肢の検討（査読待ち初稿）}\mbox{}\par

\noindent\textbf{1：不適切} 植物ホルモン／プロトプラスト／核／未受精卵。プロトプラストは細胞壁を除いた細胞であり、ここで問う不定形の細胞塊ではない。\par

\noindent\textbf{2：最適} 植物ホルモン／カルス／核／未受精卵。オーキシン等は植物ホルモンで、細胞塊はカルス。核を除いた未受精卵へ体細胞核を入れる手順にも整合。\par

\noindent\textbf{3：不適切} 植物ホルモン／カルス／胚／卵原細胞。体細胞から胚を取り出すのではなく核を使う。受容側も卵原細胞ではない。\par

\noindent\textbf{4：不適切} アミノ酸／カルス／胚／未受精卵。ア・ウが違う。植物ホルモンはアミノ酸という物質群ではなく、移す対象は核。\par

\noindent\textbf{5：不適切} アミノ酸／プロトプラスト／胚／卵原細胞。ホルモン、カルス、核、未受精卵の各対応が違う。\par

\paragraph{再利用する考え方} 単位・対象・結合の階層を区別する。植物の再分化制御と体細胞核の再プログラム化の成立条件を確認して類似年度の肢順へ対応する。

\paragraph{共通理論} \hyperref[theory:biology:cloning]{第\ref*{theory:biology:cloning}節}

\par 参考：\cite{plant_regeneration,nhgri_cloning}。

\par 難易度 2（校準初稿）、優先度未評価。

\section{材料の応力・変形}

\subsection{平成23年度 問4-4}
\label{q:h23-reg-basic-4-4}

\noindent\textbf{原典：} sources/H23/basic.pdf、PDF物理13頁。

\textbf{公式正答：} 4（single）。

\par\textit{原典本文の校正と独自解説は未作成。分類のみの登録項目。}

\subsection{平成24年度 問4-3}
\label{q:h24-reg-basic-4-3}

\noindent\textbf{原典：} sources/H24/basic.pdf、PDF物理13頁。

\textbf{公式正答：} 2（single）。

\par\textit{原典本文の校正と独自解説は未作成。分類のみの登録項目。}

\subsection{平成26年度 問4-3}
\label{q:h26-reg-basic-4-3}

\noindent\textbf{原典：} sources/H26/basic.pdf、PDF物理14頁。

\textbf{公式正答：} 5（single）。

\par\textit{原典本文の校正と独自解説は未作成。分類のみの登録項目。}

\subsection{平成28年度 問4-4}
\label{q:h28-reg-basic-4-4}

\noindent\textbf{原典：} sources/H28/basic.pdf、PDF物理18頁。

\textbf{公式正答：} 2（single）。

\par\textit{原典本文の校正と独自解説は未作成。分類のみの登録項目。}

\subsection{平成30年度 問4-4}
\label{q:h30-reg-basic-4-4}

\noindent\textbf{原典：} sources/H30/basic.pdf、PDF物理16頁。

\textbf{公式正答：} 2（single）。

\par\textit{原典本文の校正と独自解説は未作成。分類のみの登録項目。}

\subsection{令和3年度 問4-3}
\label{q:r03-reg-basic-4-3}

\noindent\textbf{原典：} sources/R03/basic.pdf、PDF物理17頁。

\textbf{公式正答：} 2（single）。

\par\textit{原典本文の校正と独自解説は未作成。分類のみの登録項目。}

\subsection{令和4年度 問4-4}
\label{q:r04-reg-basic-4-4}

\noindent\textbf{原典：} sources/R04/basic.pdf、PDF物理15頁。

\textbf{公式正答：} 1（single）。

\par\textit{原典本文の校正と独自解説は未作成。分類のみの登録項目。}

\section{生物・遺伝}

\subsection{平成23年度 問4-3}
\label{q:h23-reg-basic-4-3}

\noindent\textbf{原典：} sources/H23/basic.pdf、PDF物理12頁。

\textbf{公式正答：} 4（single）。

\par\textit{原典本文の校正と独自解説は未作成。分類のみの登録項目。}

\section{DNAの融解・再会合と紫外吸収}

\subsection{平成24年度 問4-5}
\label{q:h24-reg-basic-4-5}

\noindent\textbf{原典：} sources/H24/basic.pdf、PDF物理14頁。

\textbf{公式正答：} 3（single）。

\paragraph{問題の要旨} DNAの熱変性と再会合、吸収スペクトル、GC含量と融解温度を4空欄へ対応させる。

\Needspace{4\baselineskip}\paragraph{確認した条件・記述（独自要約）}\mbox{}\par

\noindent\textbf{空欄：} ア〜エ。熱又は強アルカリによる変性、ゆっくり冷却。\par

\noindent\textbf{比較：} GC含量と融解温度は他条件を揃える。\par

\paragraph{問われる核心} DNA鎖間会合・紫外増色と再会合

\paragraph{必要知識} DNA鎖間会合・紫外増色と再会合。熱等で二重鎖の会合が解けても通常はヌクレオチド主鎖を切断しない。変性で約260 nmの紫外吸収が増え、同じ長さ・塩濃度等ならGC含量が高いほど融解温度は一般に高い。相補鎖を適切な条件で冷却再会合させるのがアニーリングで、3。安定化には塩基スタッキングも重要であり、水素結合の本数だけを唯一の原因にしない。

\paragraph{考え方・正答を選ぶ理由（査読待ち初稿）} 熱等で二重鎖の会合が解けても通常はヌクレオチド主鎖を切断しない。変性で約260 nmの紫外吸収が増え、同じ長さ・塩濃度等ならGC含量が高いほど融解温度は一般に高い。相補鎖を適切な条件で冷却再会合させるのがアニーリングで、3。安定化には塩基スタッキングも重要であり、水素結合の本数だけを唯一の原因にしない。

\Needspace{4\baselineskip}\paragraph{選択肢の検討（査読待ち初稿）}\mbox{}\par

\noindent\textbf{1：不適切} 水素結合／紫外線／ウラシル／メルティング。鎖間会合と紫外は合うが、DNAのGの相手はCで、冷却再会合はアニーリングである。\par

\noindent\textbf{2：不適切} 共有結合／可視光線／シトシン／メルティング。通常の熱変性は主鎖の共有結合切断でなく、吸収は主に紫外であり、再会合の名称も違う。\par

\noindent\textbf{3：最適} 水素結合／紫外線／シトシン／アニーリング。鎖間会合、紫外増色、GC、冷却再会合の4つが一致する。\par

\noindent\textbf{4：不適切} 共有結合／紫外線／ウラシル／アニーリング。二重鎖を支える塩基間水素結合を主鎖共有結合と取り違え、Gの相手も違う。\par

\noindent\textbf{5：不適切} 水素結合／可視光線／シトシン／メルティング。塩基間とCは合うが、測るのは紫外、再会合はアニーリング。\par

\paragraph{再利用する考え方} 単位・対象・結合の階層を区別する。DNA鎖間会合・紫外増色と再会合の成立条件を確認して類似年度の肢順へ対応する。

\paragraph{共通理論} \hyperref[theory:biology:dna]{第\ref*{theory:biology:dna}節}, \hyperref[theory:biology:melting]{第\ref*{theory:biology:melting}節}

\par 参考：\cite{bu_dna_melting,nhgri_dna}。

\paragraph{関連問題} \hyperref[q:h28-reg-basic-4-6]{H28-REG-BASIC-4-6}

\par 難易度 2（校準初稿）、優先度未評価。

\subsection{平成28年度 問4-6}
\label{q:h28-reg-basic-4-6}

\noindent\textbf{原典：} sources/H28/basic.pdf、PDF物理20頁。

\textbf{公式正答：} 3（single）。

\paragraph{問題の要旨} DNA変性で解ける鎖間会合と残る主鎖結合、主鎖単位とGC含量を区別する。

\Needspace{4\baselineskip}\paragraph{確認した条件・記述（独自要約）}\mbox{}\par

\noindent\textbf{空欄：} ア〜エ。熱又は強アルカリで変性し、各鎖の基本構造は残る。\par

\paragraph{問われる核心} DNA鎖間と主鎖内の結合の階層

\paragraph{必要知識} DNA鎖間と主鎖内の結合の階層。塩基間の水素結合等の会合が解けても、各鎖のヌクレオチドを連ねるホスホジエステル共有結合は通常の変性で切れない。Gの相手はCである。水素・ヌクレオチド・ホスホジエステル・Cの3。融解温度はGC比だけでなく長さ・配列・塩等にも依存し、比較条件を揃える。

\paragraph{考え方・正答を選ぶ理由（査読待ち初稿）} 塩基間の水素結合等の会合が解けても、各鎖のヌクレオチドを連ねるホスホジエステル共有結合は通常の変性で切れない。Gの相手はCである。水素・ヌクレオチド・ホスホジエステル・Cの3。融解温度はGC比だけでなく長さ・配列・塩等にも依存し、比較条件を揃える。

\Needspace{4\baselineskip}\paragraph{選択肢の検討（査読待ち初稿）}\mbox{}\par

\noindent\textbf{1：不適切} ジスルフィド／ヌクレオチド／ホスホジエステル／C。主鎖単位とCは合うが、DNA塩基間はジスルフィド結合ではない。\par

\noindent\textbf{2：不適切} ジスルフィド／アミノ酸／水素／U。塩基間・核酸単位・主鎖結合・Gの相手を全て取り違える。\par

\noindent\textbf{3：最適} 水素／ヌクレオチド／ホスホジエステル／C。2鎖間の水素結合と各鎖内のヌクレオチドをつなぐ共有結合を正しく区別し、GCにも整合する。\par

\noindent\textbf{4：不適切} 水素／アミノ酸／ジスルフィド／U。鎖間は合うがアミノ酸とジスルフィドはタンパク質の概念であり、DNAのGの相手はC。\par

\noindent\textbf{5：不適切} ホスホジエステル／ヌクレオチド／ジスルフィド／C。主鎖ホスホジエステルと塩基間水素結合を逆にし、残る主鎖をジスルフィドとした点も誤り。\par

\paragraph{再利用する考え方} 単位・対象・結合の階層を区別する。DNA鎖間と主鎖内の結合の階層の成立条件を確認して類似年度の肢順へ対応する。

\paragraph{共通理論} \hyperref[theory:biology:dna]{第\ref*{theory:biology:dna}節}, \hyperref[theory:biology:melting]{第\ref*{theory:biology:melting}節}

\par 参考：\cite{bu_dna_melting,nhgri_dna}。

\paragraph{関連問題} \hyperref[q:h24-reg-basic-4-5]{H24-REG-BASIC-4-5}, \hyperref[q:r01-reg-basic-4-5]{R01-REG-BASIC-4-5}

\par 難易度 2（校準初稿）、優先度未評価。

\subsection{令和元年度 問4-5}
\label{q:r01-reg-basic-4-5}

\noindent\textbf{原典：} sources/R01/basic.pdf、PDF物理18頁。

\textbf{公式正答：} 4（single）。

\paragraph{問題の要旨} DNA変性で解ける鎖間会合と残る主鎖結合を、糖・アミノ酸等の類似語から区別する。

\Needspace{4\baselineskip}\paragraph{確認した条件・記述（独自要約）}\mbox{}\par

\noindent\textbf{空欄：} ア〜エ。候補表にグルコースが入る。\par

\paragraph{問われる核心} DNA鎖間と主鎖内の結合の階層

\paragraph{必要知識} DNA鎖間と主鎖内の結合の階層。H28と候補表が違う。原典各行を照合すると、水素・ヌクレオチド・ホスホジエステル・Cは4である。ヌクレオチドの糖はデオキシリボースで、グルコースの重合体ではない。変性と主鎖分解を区別して4を選ぶ。

\paragraph{考え方・正答を選ぶ理由（査読待ち初稿）} H28と候補表が違う。原典各行を照合すると、水素・ヌクレオチド・ホスホジエステル・Cは4である。ヌクレオチドの糖はデオキシリボースで、グルコースの重合体ではない。変性と主鎖分解を区別して4を選ぶ。

\Needspace{4\baselineskip}\paragraph{選択肢の検討（査読待ち初稿）}\mbox{}\par

\noindent\textbf{1：不適切} ジスルフィド／グルコース／水素／U。塩基間は水素、単位はヌクレオチド、主鎖はホスホジエステルであり、Gの相手はC。\par

\noindent\textbf{2：不適切} ジスルフィド／ヌクレオチド／ホスホジエステル／C。主鎖単位とCは合うが、DNA塩基間はジスルフィド結合ではない。\par

\noindent\textbf{3：不適切} 水素／グルコース／ジスルフィド／U。水素だけが合う。核酸の単位はグルコースではなくヌクレオチド、主鎖はホスホジエステル、Gの相手はC。\par

\noindent\textbf{4：最適} 水素／ヌクレオチド／ホスホジエステル／C。2鎖間の水素結合と各鎖内のヌクレオチドをつなぐ共有結合を正しく区別し、GCにも整合する。\par

\noindent\textbf{5：不適切} ホスホジエステル／ヌクレオチド／ジスルフィド／C。主鎖ホスホジエステルと塩基間水素結合を逆にし、残る主鎖をジスルフィドとした点も誤り。\par

\paragraph{再利用する考え方} 単位・対象・結合の階層を区別する。DNA鎖間と主鎖内の結合の階層の成立条件を確認して類似年度の肢順へ対応する。

\paragraph{共通理論} \hyperref[theory:biology:dna]{第\ref*{theory:biology:dna}節}, \hyperref[theory:biology:melting]{第\ref*{theory:biology:melting}節}

\par 参考：\cite{bu_dna_melting,nhgri_dna}。

\paragraph{関連問題} \hyperref[q:h28-reg-basic-4-6]{H28-REG-BASIC-4-6}

\par 難易度 2（校準初稿）、優先度未評価。

\section{DNAの相補性・主鎖・複製方向}

\subsection{平成27年度 問4-6}
\label{q:h27-reg-basic-4-6}

\noindent\textbf{原典：} sources/H27/basic.pdf、PDF物理14頁。

\textbf{公式正答：} 3（single）。

\paragraph{問題の要旨} DNAの一方の鎖のCとTの割合から、同じ鎖と相補鎖で必ず成り立つ塩基割合を求める。

\Needspace{4\baselineskip}\paragraph{確認した条件・記述（独自要約）}\mbox{}\par

\noindent\textbf{条件：} 一方の鎖C20、T35パーセント。\par

\noindent\textbf{方向：} 最も適切。割合の対象鎖を毎回確認。\par

\paragraph{問われる核心} 相補写像と一方の鎖の割合の和

\paragraph{必要知識} 相補写像と一方の鎖の割合の和。同じ鎖で\(A+G=100-(C+T)=45\)パーセント。相補鎖では\(A^{\ast}=T,\ G^{\ast}=C\)なので合計55、残り45である。3を選ぶ。二重鎖全体のA=\allowbreak{}T、G=\allowbreak{}Cを一方の鎖へ課すと1・2を誤って確定してしまう。

\paragraph{考え方・正答を選ぶ理由（査読待ち初稿）} 同じ鎖で\(A+G=100-(C+T)=45\)パーセント。相補鎖では\(A^{\ast}=T,\ G^{\ast}=C\)なので合計55、残り45である。3を選ぶ。二重鎖全体のA=\allowbreak{}T、G=\allowbreak{}Cを一方の鎖へ課すと1・2を誤って確定してしまう。

\Needspace{4\baselineskip}\paragraph{選択肢の検討（査読待ち初稿）}\mbox{}\par

\noindent\textbf{1：不適切} 同じ鎖G=\allowbreak{}20パーセント。。一方の鎖内でG=\allowbreak{}Cとは限らずG単独は未決定。\par

\noindent\textbf{2：不適切} 同じ鎖A=\allowbreak{}35パーセント。。一方の鎖内でA=\allowbreak{}Tとは限らずA単独は未決定。\par

\noindent\textbf{3：最適} 同じ鎖A+G=\allowbreak{}45パーセント。。全塩基の割合は100なので100−20−35=\allowbreak{}45。\par

\noindent\textbf{4：不適切} 相補鎖A+G=\allowbreak{}45パーセント。。相補鎖Aは元のT、Gは元のCで合計55となる。\par

\noindent\textbf{5：不適切} 相補鎖C+T=\allowbreak{}55パーセント。。相補鎖C+Tは元のG+Aなので45となる。\par

\paragraph{再利用する考え方} 単位・対象・結合の階層を区別する。相補写像と一方の鎖の割合の和の成立条件を確認して類似年度の肢順へ対応する。

\paragraph{共通理論} \hyperref[theory:biology:dna]{第\ref*{theory:biology:dna}節}, \hyperref[theory:biology:base_fraction]{第\ref*{theory:biology:base_fraction}節}

\par 参考：\cite{nhgri_dna}。

\paragraph{関連問題} \hyperref[q:r04-reg-basic-4-6]{R04-REG-BASIC-4-6}

\par 難易度 2（校準初稿）、優先度未評価。

\subsection{令和4年度 問4-6}
\label{q:r04-reg-basic-4-6}

\noindent\textbf{原典：} sources/R04/basic.pdf、PDF物理16頁。

\textbf{公式正答：} 5（single）。

\paragraph{問題の要旨} DNA一方の鎖のGとAの割合から相補鎖の和と単独割合を求める。

\Needspace{4\baselineskip}\paragraph{確認した条件・記述（独自要約）}\mbox{}\par

\noindent\textbf{条件：} 一方の鎖G25、A15パーセント。\par

\paragraph{問われる核心} 相補写像と割合の保存・非負条件

\paragraph{必要知識} 相補写像と割合の保存・非負条件。\(C+T=100-G-A=60\)パーセント。相補写像で\(G^{\ast}+A^{\ast}=C+T=60\)となり5。相補C25・T15である。一方の鎖G+Cは未決定だが、90は非負の残りを超えるため明確に排除できる。

\paragraph{考え方・正答を選ぶ理由（査読待ち初稿）} \(C+T=100-G-A=60\)パーセント。相補写像で\(G^{\ast}+A^{\ast}=C+T=60\)となり5。相補C25・T15である。一方の鎖G+Cは未決定だが、90は非負の残りを超えるため明確に排除できる。

\Needspace{4\baselineskip}\paragraph{選択肢の検討（査読待ち初稿）}\mbox{}\par

\noindent\textbf{1：不適切} 同じ鎖C+T=\allowbreak{}40パーセント。。同じ鎖の残りは100−25−15=\allowbreak{}60。\par

\noindent\textbf{2：不適切} 同じ鎖G+C=\allowbreak{}90パーセント。。G25なのでC65が必要だが、C+T60かつ非負に反する。単に未決定ではなく不可能。\par

\noindent\textbf{3：不適切} 相補鎖T=\allowbreak{}25パーセント。。相補鎖Tは元のA15に等しい。G25と対応しない。\par

\noindent\textbf{4：不適切} 相補鎖C+T=\allowbreak{}50パーセント。。相補鎖C25、T15で和40。\par

\noindent\textbf{5：最適} 相補鎖G+A=\allowbreak{}60パーセント。。相補鎖G+Aは元のC+Tなので60。\par

\paragraph{再利用する考え方} 単位・対象・結合の階層を区別する。相補写像と割合の保存・非負条件の成立条件を確認して類似年度の肢順へ対応する。

\paragraph{共通理論} \hyperref[theory:biology:base_fraction]{第\ref*{theory:biology:base_fraction}節}

\par 参考：\cite{nhgri_dna}。

\paragraph{関連問題} \hyperref[q:h27-reg-basic-4-6]{H27-REG-BASIC-4-6}

\par 難易度 2（校準初稿）、優先度未評価。

\subsection{令和6年度 問4-5}
\label{q:r06-reg-basic-4-5}

\noindent\textbf{原典：} sources/R06/basic.pdf、PDF物理16頁。

\textbf{公式正答：} 1（single）。

\paragraph{問題の要旨} DNAの塩基、逆平行と複製方向、二重らせん配置、メチル化による発現記憶を判定する。

\Needspace{4\baselineskip}\paragraph{確認した条件・記述（独自要約）}\mbox{}\par

\noindent\textbf{方向：} 最も不適切。\par

\noindent\textbf{モデル：} 二重らせん、DNAポリメラーゼによる通常の伸長。\par

\paragraph{問われる核心} DNAの構造・合成方向と発現状態の記憶

\paragraph{必要知識} DNAの構造・合成方向と発現状態の記憶。標準DNAの塩基はA,C,G,Tで、Uを含むとする1が不適切。二重らせん構造、3端への伸長、細胞分裂を通じた発現記憶を区別する。メチル化は配列変異と同一ではない。

\paragraph{考え方・正答を選ぶ理由（査読待ち初稿）} 標準DNAの塩基はA,C,G,Tで、Uを含むとする1が不適切。二重らせん構造、3端への伸長、細胞分裂を通じた発現記憶を区別する。メチル化は配列変異と同一ではない。

\Needspace{4\baselineskip}\paragraph{選択肢の検討（査読待ち初稿）}\mbox{}\par

\noindent\textbf{1：不適切} DNAの4塩基はA・C・G・U。。通常のDNAはUでなくTである。RNAの塩基と取り違える。\par

\noindent\textbf{2：適切} 逆平行2鎖を塩基間水素結合がつなぐ。。二重らせんモデルでは相補塩基間水素結合が会合へ寄与し、鎖は逆向き。スタッキングも寄与する。\par

\noindent\textbf{3：適切} DNAポリメラーゼの新生鎖は\(5'\to3'\)へ伸長する。。DNAポリメラーゼはプライマーの自由な3端へヌクレオチドを付けるので、新生鎖は5→3へ伸びる。\par

\noindent\textbf{4：適切（二重らせん）} 二重らせんで塩基は内側、糖リン酸主鎖は外側。。原典が述べる二重らせんの配置である。一本鎖DNAまで全て二重らせんであるとは読まない。\par

\noindent\textbf{5：適切（細胞分裂での機構）} 脊椎動物細胞のDNAメチル化は発現パターンを子孫細胞へ伝える機構となる。。配列自体を変えずに発現状態に関わり、分裂後の維持へ寄与するエピジェネティックな機構である。個々の印が全て無条件に維持されるとはしない。\par

\paragraph{再利用する考え方} 単位・対象・結合の階層を区別する。DNAの構造・合成方向と発現状態の記憶の成立条件を確認して類似年度の肢順へ対応する。

\paragraph{共通理論} \hyperref[theory:biology:dna]{第\ref*{theory:biology:dna}節}, \hyperref[theory:biology:epigenetics]{第\ref*{theory:biology:epigenetics}節}

\par 参考：\cite{nhgri_dna,nhgri_epigenomics}。

\par 難易度 2（校準初稿）、優先度未評価。

\section{酵素の基質特異性・障壁・温度条件}

\subsection{令和4年度 問4-5}
\label{q:r04-reg-basic-4-5}

\noindent\textbf{原典：} sources/R04/basic.pdf、PDF物理15頁。

\textbf{公式正答：} 5（single）。

\paragraph{問題の要旨} 酵素の側鎖分布、温度適応、触媒の分類と障壁、リパーゼの標的結合を判定する。

\Needspace{4\baselineskip}\paragraph{確認した条件・記述（独自要約）}\mbox{}\par

\noindent\textbf{方向：} 最も適切。\par

\paragraph{問われる核心} 酵素の基質特異性と速度・温度の条件

\paragraph{必要知識} 酵素の基質特異性と速度・温度の条件。酵素は反応速度を変える生体触媒で、一般にはタンパク質からなり、触媒RNA等もある。リパーゼがトリグリセリドのエステルを加水分解する5が適切。至適温度と安定性、反応障壁と平衡を区別する。

\paragraph{考え方・正答を選ぶ理由（査読待ち初稿）} 酵素は反応速度を変える生体触媒で、一般にはタンパク質からなり、触媒RNA等もある。リパーゼがトリグリセリドのエステルを加水分解する5が適切。至適温度と安定性、反応障壁と平衡を区別する。

\Needspace{4\baselineskip}\paragraph{選択肢の検討（査読待ち初稿）}\mbox{}\par

\noindent\textbf{1：不適切（通常の水溶性）} 非極性側鎖は酵素の外表面に多い。。通常の水溶性球状酵素では非極性Rは内部へ偏る。膜への接触面や活性部位等を区別する。\par

\noindent\textbf{2：不適切} 至適温度が20℃以下又は100℃以上の酵素は存在しない。。温度適応は多様である。実測100℃至適のPyrococcus woesei αアミラーゼだけでも不存在という断定を反証する。\par

\noindent\textbf{3：不適切} 酵素はタンパク質を主成分とする無機触媒。。多くはタンパク質、例外には触媒RNAもある生体触媒で、無機触媒という分類ではない。\par

\noindent\textbf{4：不適切} 酵素は活性化エネルギーを増加させる。。別反応経路や遷移状態の安定化で活性化障壁を下げる。反応の自由エネルギー差や平衡定数を変えるとは言わない。\par

\noindent\textbf{5：適切} リパーゼはトリグリセリドのエステル結合を加水分解する。。脂質のエステル結合に水を使って作用する酵素であり、分類も加水分解酵素。\par

\paragraph{再利用する考え方} 単位・対象・結合の階層を区別する。酵素の基質特異性と速度・温度の条件の成立条件を確認して類似年度の肢順へ対応する。

\paragraph{共通理論} \hyperref[theory:biology:enzymes]{第\ref*{theory:biology:enzymes}節}, \hyperref[theory:biology:protein]{第\ref*{theory:biology:protein}節}

\par 参考：\cite{iubmb_lipase,hyper_amylase}。

\par 難易度 2（校準初稿）、優先度未評価。

\subsection{令和5年度 問4-5}
\label{q:r05-reg-basic-4-5}

\noindent\textbf{原典：} sources/R05/basic.pdf、PDF物理14頁。

\textbf{公式正答：} 3（single）。

\paragraph{問題の要旨} タンパク質の単位・連結と生体触媒、アミラーゼの基質を4空欄へ対応させる。

\Needspace{4\baselineskip}\paragraph{確認した条件・記述（独自要約）}\mbox{}\par

\noindent\textbf{空欄：} ア〜エ。タンパク質を主成分とする生体触媒。\par

\paragraph{問われる核心} 高分子の単位・結合と加水分解の基質特異性

\paragraph{必要知識} 高分子の単位・結合と加水分解の基質特異性。タンパク質はアミノ酸のペプチド結合による鎖で、多くの酵素の主成分となる。アミラーゼはデンプンの主にα1→4結合へ作用する。セルロースはβ1→4で別の酵素を要する。4空欄が一致する3。酵素全てがタンパク質という定義へ拡張しない。

\paragraph{考え方・正答を選ぶ理由（査読待ち初稿）} タンパク質はアミノ酸のペプチド結合による鎖で、多くの酵素の主成分となる。アミラーゼはデンプンの主にα1→4結合へ作用する。セルロースはβ1→4で別の酵素を要する。4空欄が一致する3。酵素全てがタンパク質という定義へ拡張しない。

\Needspace{4\baselineskip}\paragraph{選択肢の検討（査読待ち初稿）}\mbox{}\par

\noindent\textbf{1：不適切} グルコース／イオン／酵素／デンプン。タンパク質の単位と主鎖結合が違う。グルコースは糖の単位である。\par

\noindent\textbf{2：不適切} グルコース／ペプチド／抗体／セルロース。単位が違い、抗体は一般の生体反応触媒の総称ではない。アミラーゼはセルロースのβ結合を標的としない。\par

\noindent\textbf{3：最適} アミノ酸／ペプチド／酵素／デンプン。アミノ酸、ペプチド、酵素、デンプンの4項目が整合。\par

\noindent\textbf{4：不適切} アミノ酸／ペプチド／抗体／セルロース。単位と結合は合うが、触媒の総称と基質が違う。\par

\noindent\textbf{5：不適切} アミノ酸／イオン／酵素／デンプン。主鎖は共有のペプチド結合で、イオン結合の連結ではない。\par

\paragraph{再利用する考え方} 単位・対象・結合の階層を区別する。高分子の単位・結合と加水分解の基質特異性の成立条件を確認して類似年度の肢順へ対応する。

\paragraph{共通理論} \hyperref[theory:biology:protein]{第\ref*{theory:biology:protein}節}, \hyperref[theory:biology:enzymes]{第\ref*{theory:biology:enzymes}節}

\par 参考：\cite{mit_amino,iubmb_amylase}。

\par 難易度 2（校準初稿）、優先度未評価。

\section{遺伝暗号の縮重と読み枠}

\subsection{平成28年度 問4-5}
\label{q:h28-reg-basic-4-5}

\noindent\textbf{原典：} sources/H28/basic.pdf、PDF物理19頁。

\textbf{公式正答：} 1（single）。

\paragraph{問題の要旨} 4種類の塩基の3文字列64通りと標準20アミノ酸の個数の違いを説明する。

\Needspace{4\baselineskip}\paragraph{確認した条件・記述（独自要約）}\mbox{}\par

\noindent\textbf{方向：} 64対20の説明として最も適切。\par

\noindent\textbf{範囲：} 標準遺伝暗号。特殊な再符号化・ミトコンドリアの例外を一般表へ混ぜない。\par

\paragraph{問われる核心} 遺伝暗号の多対一対応と終止

\paragraph{必要知識} 遺伝暗号の多対一対応と終止。\(4^3=64\)は3文字列の数で、出力のアミノ酸と一対一である必要はない。標準暗号では61の意味コドンが20種へ対応し、3が終止である。したがって縮重を述べる1。原典のDNAのコドンという表現と、実際にリボソームが読むmRNAのUを含むコドンを対応させる。

\paragraph{考え方・正答を選ぶ理由（査読待ち初稿）} \(4^3=64\)は3文字列の数で、出力のアミノ酸と一対一である必要はない。標準暗号では61の意味コドンが20種へ対応し、3が終止である。したがって縮重を述べる1。原典のDNAのコドンという表現と、実際にリボソームが読むmRNAのUを含むコドンを対応させる。

\Needspace{4\baselineskip}\paragraph{選択肢の検討（査読待ち初稿）}\mbox{}\par

\noindent\textbf{1：最適} 残る44コドンの多くも同じ20アミノ酸へ対応し、1種に複数コドンがある。。標準遺伝暗号の縮重であり、多対一の対応を表す。61がアミノ酸、3が終止なので44全てをアミノ酸へ割り当てたとは言わない。\par

\noindent\textbf{2：不適切} 残る44コドンはアミノ酸の修飾体へ用いられる。。翻訳後の修飾等を専用の44コドンで指定する仕組みではない。\par

\noindent\textbf{3：不適切} 転写で配列が変わり1アミノ酸1コドンへ収束する。。通常の転写は相補配列を作り、同義コドンを全て1個へ書き換える過程ではない。\par

\noindent\textbf{4：不適切} 進化で残る44コドンが遺伝子から消えた。。同義コドンは現在の遺伝子でも使用される。存在しないという説明は違う。\par

\noindent\textbf{5：不適切} コドンの第1塩基は通常遺伝情報として意味を持たない。。第1位置の変化はアミノ酸を変え得る。第3位置の縮重も全コドンで無意味ということではない。\par

\paragraph{再利用する考え方} 単位・対象・結合の階層を区別する。遺伝暗号の多対一対応と終止の成立条件を確認して類似年度の肢順へ対応する。

\paragraph{共通理論} \hyperref[theory:biology:code]{第\ref*{theory:biology:code}節}

\par 参考：\cite{nhgri_codon}。

\paragraph{関連問題} \hyperref[q:r01-ret-basic-4-5]{R01-RET-BASIC-4-5}

\par 難易度 2（校準初稿）、優先度未評価。

\subsection{令和元年度再試験 問4-5}
\label{q:r01-ret-basic-4-5}

\noindent\textbf{原典：} sources/R01\_retest/basic.pdf、PDF物理15頁。

\textbf{公式正答：} 3（single）。

\paragraph{問題の要旨} 4種類の塩基の3文字列64通りと標準20アミノ酸の個数の違いを説明する。

\Needspace{4\baselineskip}\paragraph{確認した条件・記述（独自要約）}\mbox{}\par

\noindent\textbf{方向：} 64対20の説明として最も適切。\par

\noindent\textbf{範囲：} 標準遺伝暗号。特殊な再符号化・ミトコンドリアの例外を一般表へ混ぜない。\par

\paragraph{問われる核心} 遺伝暗号の多対一対応と終止

\paragraph{必要知識} 遺伝暗号の多対一対応と終止。\(4^3=64\)は3文字列の数で、出力のアミノ酸と一対一である必要はない。標準暗号では61の意味コドンが20種へ対応し、3が終止である。したがって縮重を述べる3。原典のDNAのコドンという表現と、実際にリボソームが読むmRNAのUを含むコドンを対応させる。

\paragraph{考え方・正答を選ぶ理由（査読待ち初稿）} \(4^3=64\)は3文字列の数で、出力のアミノ酸と一対一である必要はない。標準暗号では61の意味コドンが20種へ対応し、3が終止である。したがって縮重を述べる3。原典のDNAのコドンという表現と、実際にリボソームが読むmRNAのUを含むコドンを対応させる。

\Needspace{4\baselineskip}\paragraph{選択肢の検討（査読待ち初稿）}\mbox{}\par

\noindent\textbf{1：不適切} コドンの第1塩基は通常遺伝情報として意味を持たない。。第1位置の変化はアミノ酸を変え得る。第3位置の縮重も全コドンで無意味ということではない。\par

\noindent\textbf{2：不適切} 進化で残る44コドンが遺伝子から消えた。。同義コドンは現在の遺伝子でも使用される。存在しないという説明は違う。\par

\noindent\textbf{3：最適} 残る44コドンの多くも同じ20アミノ酸へ対応し、1種に複数コドンがある。。標準遺伝暗号の縮重であり、多対一の対応を表す。61がアミノ酸、3が終止なので44全てをアミノ酸へ割り当てたとは言わない。\par

\noindent\textbf{4：不適切} 転写で配列が変わり1アミノ酸1コドンへ収束する。。通常の転写は相補配列を作り、同義コドンを全て1個へ書き換える過程ではない。\par

\noindent\textbf{5：不適切} 残る44コドンはアミノ酸の修飾体へ用いられる。。翻訳後の修飾等を専用の44コドンで指定する仕組みではない。\par

\paragraph{再利用する考え方} 単位・対象・結合の階層を区別する。遺伝暗号の多対一対応と終止の成立条件を確認して類似年度の肢順へ対応する。

\paragraph{共通理論} \hyperref[theory:biology:code]{第\ref*{theory:biology:code}節}

\par 参考：\cite{nhgri_codon}。

\paragraph{関連問題} \hyperref[q:h28-reg-basic-4-5]{H28-REG-BASIC-4-5}

\par 難易度 2（校準初稿）、優先度未評価。

\section{塩基変化の種類と遺伝子発現への効果}

\subsection{令和3年度 問4-6}
\label{q:r03-reg-basic-4-6}

\noindent\textbf{原典：} sources/R03/basic.pdf、PDF物理20頁。

\textbf{公式正答：} 2（single）。

\paragraph{問題の要旨} 塩基置換・挿入・欠失の翻訳への効果と潜性形質、変異誘発を判定する。

\Needspace{4\baselineskip}\paragraph{確認した条件・記述（独自要約）}\mbox{}\par

\noindent\textbf{方向：} 最も適切。\par

\paragraph{問われる核心} 読み枠の変化と塩基置換の効果の区別

\paragraph{必要知識} 読み枠の変化と塩基置換の効果の区別。読み枠は3塩基単位なので、翻訳領域で3の倍数でない挿入・欠失は下流の区切りを変える。1塩基挿入の影響を述べる2が適切。他の肢は変異の名称・分子変化の種別・潜性のモデル又は変異原の作用を取り違える。

\paragraph{考え方・正答を選ぶ理由（査読待ち初稿）} 読み枠は3塩基単位なので、翻訳領域で3の倍数でない挿入・欠失は下流の区切りを変える。1塩基挿入の影響を述べる2が適切。他の肢は変異の名称・分子変化の種別・潜性のモデル又は変異原の作用を取り違える。

\Needspace{4\baselineskip}\paragraph{選択肢の検討（査読待ち初稿）}\mbox{}\par

\noindent\textbf{1：不適切} 途中の終止コドンへの置換を中立突然変異という。。早期終止はナンセンス変異であり、中立という呼び方は影響・選択上の別軸である。\par

\noindent\textbf{2：最適} 1塩基挿入で読み枠がずれアミノ酸配列が大きく変わる可能性が高い。。通常の翻訳領域に1塩基が入り、3の倍数でないため以降の読み枠がずれる。全ての非翻訳領域の挿入まで同じ効果とはしない。\par

\noindent\textbf{3：不適切} 鎌状赤血球のGlu→Valは1塩基欠失による。。代表的なHbSはHBBの塩基置換で、欠失によるフレームずれとは違う。\par

\noindent\textbf{4：不適切} 潜性形質の発現は相同染色体の一方の変異だけでよい。。常染色体の通常の二倍体・完全潜性というモデルでは、正常な対立遺伝子が一方にあれば形質は隠れ、両方の影響を考える。半接合等は別条件。\par

\noindent\textbf{5：不適切} X線・UV・化学物質で遺伝子変異は起こりにくい。。放射線・UVや特定の化学物質はDNA損傷・変異の誘発源となる。全ての化学物質が一律に変異原という意味ではない。\par

\paragraph{再利用する考え方} 単位・対象・結合の階層を区別する。読み枠の変化と塩基置換の効果の区別の成立条件を確認して類似年度の肢順へ対応する。

\paragraph{共通理論} \hyperref[theory:biology:code]{第\ref*{theory:biology:code}節}, \hyperref[theory:biology:mutation]{第\ref*{theory:biology:mutation}節}

\par 参考：\cite{nhgri_frameshift,nhgri_sickle}。

\par 難易度 2（校準初稿）、優先度未評価。

\section{同位体の基礎}

\subsection{平成28年度 問4-2}
\label{q:h28-reg-basic-4-2}

\noindent\textbf{原典：} sources/H28/basic.pdf、PDF物理17頁。

\textbf{公式正答：} 4（single）。

\paragraph{問題の要旨} 中性のCa、Ar、Cl核種の比較について、適切な記述を選ぶ。

\Needspace{4\baselineskip}\paragraph{確認した条件・記述（独自要約）}\mbox{}\par

\noindent\textbf{核種：} \({}^{40}_{20}Ca,{}^{40}_{18}Ar,{}^{35}_{17}Cl,{}^{37}_{17}Cl\)\par

\noindent\textbf{条件：} いずれも電荷を持たない原子。\par

\paragraph{問われる核心} 核種のZ・Aと中性原子の電子数

\paragraph{必要知識} 質量数は核子数、N=\allowbreak{}A−Z、電子数=Z、同位体・同重体・同素体。

\paragraph{考え方・正答を選ぶ理由（査読待ち初稿）} 陽子数Z、質量数Aより中性子数はA−Z、中性原子の電子数はZである。Ca-40は20、Ar-40は22中性子、Cl-35は18、Cl-37は20中性子となる。両Clの電子は17で等しく4が正しい。イオンなら電子数は電荷だけ変わるが、この問題は中性を指定している。

\Needspace{4\baselineskip}\paragraph{選択肢の検討（査読待ち初稿）}\mbox{}\par

\noindent\textbf{1：不適切} Ca-40(Z=\allowbreak{}20)とAr-40(Z=\allowbreak{}18)の中性子数は等しい。。中性子数A−Zは20と22なので違う。\par

\noindent\textbf{2：不適切} Cl-35とCl-37(Z=\allowbreak{}17)の中性子数は等しい。。35−17=\allowbreak{}18、37−17=\allowbreak{}20なので違う。\par

\noindent\textbf{3：不適切} Ca-40とAr-40は同位体である。。同位体はZが同じでAが違う。ここではZが違いAが同じ同重体である。\par

\noindent\textbf{4：適切} Cl-35とCl-37の電子数は等しい。。中性原子なので電子数はいずれもZ=\allowbreak{}17である。\par

\noindent\textbf{5：不適切} Cl-35とCl-37は同素体である。。同素体は同じ元素が異なる結合・構造を持つ単体であり、核の中性子数違いとは異なる。\par

\paragraph{再利用する考え方} 核種のZ・Aと中性原子の電子数を、数える対象・分母・電荷保存・成立条件を明示して判断する。

\paragraph{共通理論} \hyperref[theory:basicchem:isotope]{第\ref*{theory:basicchem:isotope}節}

\par 参考：\cite{iupac_isotopes}。

\paragraph{関連問題} \hyperref[q:r05-reg-basic-4-1]{R05-REG-BASIC-4-1}

\par 難易度 2（校準初稿）、優先度未評価。

\subsection{令和元年度 問4-2}
\label{q:r01-reg-basic-4-2}

\noindent\textbf{原典：} sources/R01/basic.pdf、PDF物理16頁。

\textbf{公式正答：} 5（single）。

\paragraph{問題の要旨} 同位体について5記述の正誤を判定し、その組合せを選ぶ。

\Needspace{4\baselineskip}\paragraph{確認した条件・記述（独自要約）}\mbox{}\par

\noindent\textbf{ア：} ア：陽子数は同じだが電子数は違う。\par

\noindent\textbf{イ：} イ：質量数が違うため化学的性質も違う。\par

\noindent\textbf{ウ：} ウ：中性子数が違う。\par

\noindent\textbf{エ：} エ：放射性同位体は医療や考古学に使う。\par

\noindent\textbf{オ：} オ：放射性同位体には別元素へ変化するものがある。\par

\noindent\textbf{並び：} ア／イ／ウ／エ／オ。一般の中性原子を基礎化学の近似で比較する。\par

\paragraph{問われる核心} 核の違い・電子配置・放射壊変を区別する

\paragraph{必要知識} 同位体定義、同位体効果、核種変化、放射性年代測定・診断。

\paragraph{考え方・正答を選ぶ理由（査読待ち初稿）} 同位体は同じZで異なるAを持ち、中性子数が違う。中性なら電子数も電子配置も同じなので基礎的化学性質はほぼ同じで、ア・イを誤、ウを正とする。放射性核種は医療・年代測定に利用でき、α崩壊やβ崩壊でZが変わる場合があるためエ・オは正。誤／誤／正／正／正の5である。ただし「同じ化学的性質」は同位体効果まで0という意味ではなく、速度や平衡に質量由来の差が現れる。全放射壊変で元素が変わるとも限らず、γ遷移はZを変えない。

\Needspace{4\baselineskip}\paragraph{選択肢の検討（査読待ち初稿）}\mbox{}\par

\noindent\textbf{1：条件に不一致} 正／正／誤／誤／誤。ア・イ・ウ・エ・オの正誤が違う。同じ電子配置、中性子差、放射性核種の応用と元素変化で判定する。\par

\noindent\textbf{2：条件に不一致} 正／正／正／正／誤。ア・イ・オの正誤が違う。同じ電子配置、中性子差、放射性核種の応用と元素変化で判定する。\par

\noindent\textbf{3：条件に不一致} 誤／誤／正／誤／誤。エ・オの正誤が違う。同じ電子配置、中性子差、放射性核種の応用と元素変化で判定する。\par

\noindent\textbf{4：条件に不一致} 誤／正／誤／正／正。イ・ウの正誤が違う。同じ電子配置、中性子差、放射性核種の応用と元素変化で判定する。\par

\noindent\textbf{5：条件に一致} 誤／誤／正／正／正。全5項の判断に一致する。\par

\paragraph{再利用する考え方} 核の違い・電子配置・放射壊変を区別するを、数える対象・分母・電荷保存・成立条件を明示して判断する。

\paragraph{共通理論} \hyperref[theory:basicchem:isotope]{第\ref*{theory:basicchem:isotope}節}, \hyperref[theory:basicchem:isotope_effect]{第\ref*{theory:basicchem:isotope_effect}節}

\par 参考：\cite{iupac_isotopes,iupac_isotope_effect}。

\paragraph{関連問題} \hyperref[q:r03-reg-basic-4-1]{R03-REG-BASIC-4-1}

\par 難易度 2（校準初稿）、優先度未評価。

\subsection{令和3年度 問4-1}
\label{q:r03-reg-basic-4-1}

\noindent\textbf{原典：} sources/R03/basic.pdf、PDF物理16頁。

\textbf{公式正答：} 3（single）。

\paragraph{問題の要旨} 同位体について5記述の正誤を判定し、その組合せを選ぶ。

\Needspace{4\baselineskip}\paragraph{確認した条件・記述（独自要約）}\mbox{}\par

\noindent\textbf{ア：} ア：質量数が違うため化学的性質も違う。\par

\noindent\textbf{イ：} イ：陽子数は同じだが電子数は違う。\par

\noindent\textbf{ウ：} ウ：中性子数が違う。\par

\noindent\textbf{エ：} エ：放射性同位体には別元素へ変化するものがある。\par

\noindent\textbf{オ：} オ：放射性同位体は年代測定等へ利用する。\par

\noindent\textbf{並び：} ア／イ／ウ／エ／オ。一般の中性原子を基礎化学の近似で比較する。\par

\paragraph{問われる核心} 核の違い・電子配置・放射壊変を区別する

\paragraph{必要知識} 同位体定義、同位体効果、核種変化、放射性年代測定・診断。

\paragraph{考え方・正答を選ぶ理由（査読待ち初稿）} 同位体は同じZで異なるAを持ち、中性子数が違う。中性なら電子数も電子配置も同じなので基礎的化学性質はほぼ同じで、ア・イを誤、ウを正とする。放射性核種は医療・年代測定に利用でき、α崩壊やβ崩壊でZが変わる場合があるためエ・オは正。誤／誤／正／正／正の3である。ただし「同じ化学的性質」は同位体効果まで0という意味ではなく、速度や平衡に質量由来の差が現れる。全放射壊変で元素が変わるとも限らず、γ遷移はZを変えない。

\Needspace{4\baselineskip}\paragraph{選択肢の検討（査読待ち初稿）}\mbox{}\par

\noindent\textbf{1：条件に不一致} 正／正／誤／誤／誤。ア・イ・ウ・エ・オの正誤が違う。同じ電子配置、中性子差、放射性核種の応用と元素変化で判定する。\par

\noindent\textbf{2：条件に不一致} 正／正／正／正／誤。ア・イ・オの正誤が違う。同じ電子配置、中性子差、放射性核種の応用と元素変化で判定する。\par

\noindent\textbf{3：条件に一致} 誤／誤／正／正／正。全5項の判断に一致する。\par

\noindent\textbf{4：条件に不一致} 誤／正／誤／正／正。イ・ウの正誤が違う。同じ電子配置、中性子差、放射性核種の応用と元素変化で判定する。\par

\noindent\textbf{5：条件に不一致} 誤／誤／正／誤／誤。エ・オの正誤が違う。同じ電子配置、中性子差、放射性核種の応用と元素変化で判定する。\par

\paragraph{再利用する考え方} 核の違い・電子配置・放射壊変を区別するを、数える対象・分母・電荷保存・成立条件を明示して判断する。

\paragraph{共通理論} \hyperref[theory:basicchem:isotope]{第\ref*{theory:basicchem:isotope}節}, \hyperref[theory:basicchem:isotope_effect]{第\ref*{theory:basicchem:isotope_effect}節}

\par 参考：\cite{iupac_isotopes,iupac_isotope_effect}。

\paragraph{関連問題} \hyperref[q:r01-reg-basic-4-2]{R01-REG-BASIC-4-2}

\par 難易度 2（校準初稿）、優先度未評価。

\subsection{令和5年度 問4-1}
\label{q:r05-reg-basic-4-1}

\noindent\textbf{原典：} sources/R05/basic.pdf、PDF物理13頁。

\textbf{公式正答：} 3（single）。

\paragraph{問題の要旨} 中性のCa、Ar、Cl核種の比較について、適切な記述を選ぶ。

\Needspace{4\baselineskip}\paragraph{確認した条件・記述（独自要約）}\mbox{}\par

\noindent\textbf{核種：} \({}^{40}_{20}Ca,{}^{40}_{18}Ar,{}^{35}_{17}Cl,{}^{37}_{17}Cl\)\par

\noindent\textbf{条件：} いずれも電荷を持たない原子。\par

\paragraph{問われる核心} 核種のZ・Aと中性原子の電子数

\paragraph{必要知識} 質量数は核子数、N=\allowbreak{}A−Z、電子数=Z、同位体・同重体・同素体。

\paragraph{考え方・正答を選ぶ理由（査読待ち初稿）} 陽子数Z、質量数Aより中性子数はA−Z、中性原子の電子数はZである。Ca-40は20、Ar-40は22中性子、Cl-35は18、Cl-37は20中性子となる。両Clの電子は17で等しく3が正しい。イオンなら電子数は電荷だけ変わるが、この問題は中性を指定している。

\Needspace{4\baselineskip}\paragraph{選択肢の検討（査読待ち初稿）}\mbox{}\par

\noindent\textbf{1：不適切} Ca-40(Z=\allowbreak{}20)とAr-40(Z=\allowbreak{}18)の中性子数は等しい。。中性子数A−Zは20と22なので違う。\par

\noindent\textbf{2：不適切} Cl-35とCl-37(Z=\allowbreak{}17)の中性子数は等しい。。35−17=\allowbreak{}18、37−17=\allowbreak{}20なので違う。\par

\noindent\textbf{3：適切} Cl-35とCl-37の電子数は等しい。。中性原子なので電子数はいずれもZ=\allowbreak{}17である。\par

\noindent\textbf{4：不適切} Ca-40とAr-40は同位体である。。同位体はZが同じでAが違う。ここではZが違いAが同じ同重体である。\par

\noindent\textbf{5：不適切} Cl-35とCl-37は同素体である。。同素体は同じ元素が異なる結合・構造を持つ単体であり、核の中性子数違いとは異なる。\par

\paragraph{再利用する考え方} 核種のZ・Aと中性原子の電子数を、数える対象・分母・電荷保存・成立条件を明示して判断する。

\paragraph{共通理論} \hyperref[theory:basicchem:isotope]{第\ref*{theory:basicchem:isotope}節}

\par 参考：\cite{iupac_isotopes}。

\paragraph{関連問題} \hyperref[q:h28-reg-basic-4-2]{H28-REG-BASIC-4-2}

\par 難易度 2（校準初稿）、優先度未評価。

\section{膜脂質の充填と流動性}

\subsection{平成23年度 問4-5}
\label{q:h23-reg-basic-4-5}

\noindent\textbf{原典：} sources/H23/basic.pdf、PDF物理14頁。

\textbf{公式正答：} 3（single）。

\paragraph{問題の要旨} 膜脂質の脂肪酸の結合形式、鎖長・不飽和化による融点と流動性、昇温への組成適応を6空欄へ対応させる。

\Needspace{4\baselineskip}\paragraph{確認した条件・記述（独自要約）}\mbox{}\par

\noindent\textbf{空欄：} ア・イ・ウ・エ・オ・カ。\par

\noindent\textbf{温度：} 細菌培養20→30℃。増える場合があるという条件。\par

\noindent\textbf{比較：} 12炭素以上の飽和鎖長と同炭素数の飽和/不飽和。通常のcis型で説明。\par

\paragraph{問われる核心} 両親媒性脂質の充填と恒常的な膜流動性

\paragraph{必要知識} 両親媒性脂質の充填と恒常的な膜流動性。通常のグリセロリン脂質では脂肪酸はエステルとして結合する。長い飽和鎖は接触と相互作用が増して融点が高く、cis二重結合は密な充填を妨げ、同炭素数の飽和鎖より融点を下げ流動性を増す。20℃から30℃への昇温で流動化するため、長鎖・飽和脂肪酸の増加で相殺する場合がある。順にエステル・高い・長鎖・低い・増す・飽和で3。全ての脂質をエステルとせず、スフィンゴ脂質のアミドや古細菌のエーテルを区別する。

\paragraph{考え方・正答を選ぶ理由（査読待ち初稿）} 通常のグリセロリン脂質では脂肪酸はエステルとして結合する。長い飽和鎖は接触と相互作用が増して融点が高く、cis二重結合は密な充填を妨げ、同炭素数の飽和鎖より融点を下げ流動性を増す。20℃から30℃への昇温で流動化するため、長鎖・飽和脂肪酸の増加で相殺する場合がある。順にエステル・高い・長鎖・低い・増す・飽和で3。全ての脂質をエステルとせず、スフィンゴ脂質のアミドや古細菌のエーテルを区別する。

\Needspace{4\baselineskip}\paragraph{選択肢の検討（査読待ち初稿）}\mbox{}\par

\noindent\textbf{1：不適切} アミド／低い／長鎖脂肪酸／高い／減る／不飽和脂肪酸。通常のグリセロリン脂質の脂肪酸はエステル結合である。長い飽和鎖の融点と不飽和化の流動性も逆。\par

\noindent\textbf{2：不適切} アミド／高い／短鎖脂肪酸／低い／増す／飽和脂肪酸。融点・不飽和化の方向は合うが、主題のエステルと、昇温に応じた長鎖化を満たさない。\par

\noindent\textbf{3：最適} エステル／高い／長鎖脂肪酸／低い／増す／飽和脂肪酸。6空欄の全てが通常の膜脂質と温度適応に整合する。\par

\noindent\textbf{4：不適切} エステル／低い／長鎖脂肪酸／高い／増す／不飽和脂肪酸。飽和鎖長と融点、不飽和化と融点、昇温への補償の方向が違う。\par

\noindent\textbf{5：不適切} エステル／高い／短鎖脂肪酸／低い／減る／不飽和脂肪酸。昇温で過度な流動化を抑えるなら短鎖化・不飽和化でなく、その逆が候補である。不飽和化で流動性が減る点も違う。\par

\paragraph{再利用する考え方} 単位・対象・結合の階層を区別する。両親媒性脂質の充填と恒常的な膜流動性の成立条件を確認して類似年度の肢順へ対応する。

\paragraph{共通理論} \hyperref[theory:biology:membrane]{第\ref*{theory:biology:membrane}節}

\par 参考：\cite{bio_membrane}。

\paragraph{関連問題} \hyperref[q:h26-reg-basic-4-5]{H26-REG-BASIC-4-5}

\par 難易度 2（校準初稿）、優先度未評価。

\subsection{平成26年度 問4-5}
\label{q:h26-reg-basic-4-5}

\noindent\textbf{原典：} sources/H26/basic.pdf、PDF物理15頁。

\textbf{公式正答：} 3（single）。

\paragraph{問題の要旨} 膜脂質の脂肪酸の結合形式、鎖長・不飽和化による融点と流動性、昇温への組成適応を6空欄へ対応させる。

\Needspace{4\baselineskip}\paragraph{確認した条件・記述（独自要約）}\mbox{}\par

\noindent\textbf{空欄：} ア・イ・ウ・エ・オ・カ。\par

\noindent\textbf{温度：} 細菌培養20→30℃。増える場合があるという条件。\par

\noindent\textbf{比較：} 12炭素以上の飽和鎖長と同炭素数の飽和/不飽和。通常のcis型で説明。\par

\paragraph{問われる核心} 両親媒性脂質の充填と恒常的な膜流動性

\paragraph{必要知識} 両親媒性脂質の充填と恒常的な膜流動性。通常のグリセロリン脂質では脂肪酸はエステルとして結合する。長い飽和鎖は接触と相互作用が増して融点が高く、cis二重結合は密な充填を妨げ、同炭素数の飽和鎖より融点を下げ流動性を増す。20℃から30℃への昇温で流動化するため、長鎖・飽和脂肪酸の増加で相殺する場合がある。順にエステル・高い・長鎖・低い・増す・飽和で3。全ての脂質をエステルとせず、スフィンゴ脂質のアミドや古細菌のエーテルを区別する。

\paragraph{考え方・正答を選ぶ理由（査読待ち初稿）} 通常のグリセロリン脂質では脂肪酸はエステルとして結合する。長い飽和鎖は接触と相互作用が増して融点が高く、cis二重結合は密な充填を妨げ、同炭素数の飽和鎖より融点を下げ流動性を増す。20℃から30℃への昇温で流動化するため、長鎖・飽和脂肪酸の増加で相殺する場合がある。順にエステル・高い・長鎖・低い・増す・飽和で3。全ての脂質をエステルとせず、スフィンゴ脂質のアミドや古細菌のエーテルを区別する。

\Needspace{4\baselineskip}\paragraph{選択肢の検討（査読待ち初稿）}\mbox{}\par

\noindent\textbf{1：不適切} アミド／低い／長鎖脂肪酸／高い／減る／不飽和脂肪酸。通常のグリセロリン脂質の脂肪酸はエステル結合である。長い飽和鎖の融点と不飽和化の流動性も逆。\par

\noindent\textbf{2：不適切} アミド／高い／短鎖脂肪酸／低い／増す／飽和脂肪酸。融点・不飽和化の方向は合うが、主題のエステルと、昇温に応じた長鎖化を満たさない。\par

\noindent\textbf{3：最適} エステル／高い／長鎖脂肪酸／低い／増す／飽和脂肪酸。6空欄の全てが通常の膜脂質と温度適応に整合する。\par

\noindent\textbf{4：不適切} エステル／低い／長鎖脂肪酸／高い／増す／不飽和脂肪酸。飽和鎖長と融点、不飽和化と融点、昇温への補償の方向が違う。\par

\noindent\textbf{5：不適切} エステル／高い／短鎖脂肪酸／低い／減る／不飽和脂肪酸。昇温で過度な流動化を抑えるなら短鎖化・不飽和化でなく、その逆が候補である。不飽和化で流動性が減る点も違う。\par

\paragraph{再利用する考え方} 単位・対象・結合の階層を区別する。両親媒性脂質の充填と恒常的な膜流動性の成立条件を確認して類似年度の肢順へ対応する。

\paragraph{共通理論} \hyperref[theory:biology:membrane]{第\ref*{theory:biology:membrane}節}

\par 参考：\cite{bio_membrane}。

\paragraph{関連問題} \hyperref[q:h23-reg-basic-4-5]{H23-REG-BASIC-4-5}

\par 難易度 2（校準初稿）、優先度未評価。

\section{PCR・制限酵素・電気泳動・DNAライブラリ}

\subsection{平成29年度 問4-6}
\label{q:h29-reg-basic-4-6}

\noindent\textbf{原典：} sources/H29/basic.pdf、PDF物理17頁。

\textbf{公式正答：} 1（single）。

\paragraph{問題の要旨} PCRの酵素、組換えインスリンの時点、電気泳動、制限酵素、ゲノムとcDNAを判定する。

\Needspace{4\baselineskip}\paragraph{確認した条件・記述（独自要約）}\mbox{}\par

\noindent\textbf{方向：} 最も適切。\par

\paragraph{問われる核心} 分子生物学の操作対象・電荷・確率モデル

\paragraph{必要知識} 分子生物学の操作対象・電荷・確率モデル。熱変性を含む通常のPCRで耐熱性ポリメラーゼを用いる1が適切。2は歴史、3は電荷の符号、4は期待間隔と個別の実現値、5はゲノム全体と組織の発現RNAを取り違える。

\paragraph{考え方・正答を選ぶ理由（査読待ち初稿）} 熱変性を含む通常のPCRで耐熱性ポリメラーゼを用いる1が適切。2は歴史、3は電荷の符号、4は期待間隔と個別の実現値、5はゲノム全体と組織の発現RNAを取り違える。

\Needspace{4\baselineskip}\paragraph{選択肢の検討（査読待ち初稿）}\mbox{}\par

\noindent\textbf{1：最適} PCRの反復熱変性のため耐熱性DNAポリメラーゼを使う。。温度サイクルの変性段階にも耐える酵素を利用し、各回に酵素を加え直す必要を避ける。\par

\noindent\textbf{2：不適切} 大腸菌による組換えインスリン合成の成功は1990年代後半。。組換えヒトインスリンは1978年の製法開発、1982年の米承認であり、1990年代後半ではない。\par

\noindent\textbf{3：不適切} DNAはゲル電気泳動で陰極へ移動する。。通常の緩衝液でDNAのリン酸骨格は負に帯電するため陽極へ進む。\par

\noindent\textbf{4：不適切} EcoRIの6塩基認識で切断後は全て正確に\(4^6\)塩基対。。独立・等確率の4塩基というランダムモデルで認識配列の出現確率は1/4096にすぎず、各断片の長さを固定しない。実ゲノムの塩基偏り・配列相関も別である。\par

\noindent\textbf{5：不適切} ヒトゲノムの全クローンが肝臓RNA由来cDNAライブラリにある。。ゲノムは非転写領域・イントロン等も含む。cDNAは採取RNAから作り、その組織・時点で発現していない遺伝子も欠け得る。\par

\paragraph{再利用する考え方} 単位・対象・結合の階層を区別する。分子生物学の操作対象・電荷・確率モデルの成立条件を確認して類似年度の肢順へ対応する。

\paragraph{共通理論} \hyperref[theory:biology:pcr]{第\ref*{theory:biology:pcr}節}, \hyperref[theory:biology:libraries]{第\ref*{theory:biology:libraries}節}

\par 参考：\cite{neb_pcr,fda_insulin,neb_ecori,nhgri_cdna}。

\paragraph{関連問題} \hyperref[q:r01-ret-basic-4-6]{R01-RET-BASIC-4-6}

\par 難易度 2（校準初稿）、優先度未評価。

\subsection{令和元年度再試験 問4-6}
\label{q:r01-ret-basic-4-6}

\noindent\textbf{原典：} sources/R01\_retest/basic.pdf、PDF物理16頁。

\textbf{公式正答：} 3（single）。

\paragraph{問題の要旨} 組換えDNA技術での発現、PCR増幅率、制限酵素の期待間隔、電気泳動方向を判定する。

\Needspace{4\baselineskip}\paragraph{確認した条件・記述（独自要約）}\mbox{}\par

\noindent\textbf{方向：} 最も適切。\par

\paragraph{問われる核心} ゲノムの存在と組織・時点別の発現の区別

\paragraph{必要知識} ゲノムの存在と組織・時点別の発現の区別。ゲノムライブラリとcDNAライブラリは由来が違う。DNAにある翻訳領域でも、その組織・時点でRNAになっていなければcDNAに入らず、3。理想増幅は\(N_n=N_0\,2^n\)で、30回では約\(1.07\times10^9\)倍である。

\paragraph{考え方・正答を選ぶ理由（査読待ち初稿）} ゲノムライブラリとcDNAライブラリは由来が違う。DNAにある翻訳領域でも、その組織・時点でRNAになっていなければcDNAに入らず、3。理想増幅は\(N_n=N_0\,2^n\)で、30回では約\(1.07\times10^9\)倍である。

\Needspace{4\baselineskip}\paragraph{選択肢の検討（査読待ち初稿）}\mbox{}\par

\noindent\textbf{1：不適切} 大腸菌による組換えインスリン合成の成功は1990年代後半。。組換えヒトインスリンは1978年の製法開発、1982年の米承認であり、1990年代後半ではない。\par

\noindent\textbf{2：不適切} 理想的PCRで30回反復すると鋳型DNAは30倍。。完全な増幅が各回続く理想なら指数増幅で2の30乗。実反応の効率・枯渇を含む定数倍30ではない。\par

\noindent\textbf{3：最適} 翻訳領域が同じ組織のゲノムライブラリにあってもcDNAにない場合がある。。その組織・採取時点で当該遺伝子が発現していなければRNA由来cDNAにはない。ゲノムに領域が存在することと転写は別。\par

\noindent\textbf{4：不適切} EcoRIの6塩基認識で切断後は全て正確に\(4^6\)塩基対。。独立・等確率の4塩基というランダムモデルで認識配列の出現確率は1/4096にすぎず、各断片の長さを固定しない。実ゲノムの塩基偏り・配列相関も別である。\par

\noindent\textbf{5：不適切} DNAはゲル電気泳動で陰極へ移動する。。通常の緩衝液でDNAのリン酸骨格は負に帯電するため陽極へ進む。\par

\paragraph{再利用する考え方} 単位・対象・結合の階層を区別する。ゲノムの存在と組織・時点別の発現の区別の成立条件を確認して類似年度の肢順へ対応する。

\paragraph{共通理論} \hyperref[theory:biology:pcr]{第\ref*{theory:biology:pcr}節}, \hyperref[theory:biology:libraries]{第\ref*{theory:biology:libraries}節}

\par 参考：\cite{neb_pcr,fda_insulin,neb_ecori,nhgri_cdna}。

\paragraph{関連問題} \hyperref[q:h29-reg-basic-4-6]{H29-REG-BASIC-4-6}

\par 難易度 2（校準初稿）、優先度未評価。

\subsection{令和2年度 問4-6}
\label{q:r02-reg-basic-4-6}

\noindent\textbf{原典：} sources/R02/basic.pdf、PDF物理18頁。

\textbf{公式正答：} 4（single）。

\paragraph{問題の要旨} PCRの熱変性・アニーリング・伸長の条件と産物に残るプライマーを判定する。

\Needspace{4\baselineskip}\paragraph{確認した条件・記述（独自要約）}\mbox{}\par

\noindent\textbf{方向：} 最も適切。\par

\paragraph{問われる核心} PCRの3段階と各段階の制御量

\paragraph{必要知識} PCRの3段階と各段階の制御量。DNA鎖を分け、標的へプライマーを会合させ、その3端から合成する。温度サイクルに耐える酵素を用いる4が適切。伸長時間は配列長と酵素速度に応じて確保し、温度と特異性の関係を反転しない。

\paragraph{考え方・正答を選ぶ理由（査読待ち初稿）} DNA鎖を分け、標的へプライマーを会合させ、その3端から合成する。温度サイクルに耐える酵素を用いる4が適切。伸長時間は配列長と酵素速度に応じて確保し、温度と特異性の関係を反転しない。

\Needspace{4\baselineskip}\paragraph{選択肢の検討（査読待ち初稿）}\mbox{}\par

\noindent\textbf{1：不適切} 熱変性はDNA主鎖の共有結合を切断して一本鎖へ解離させる。。通常の変性は鎖間の会合を解く段階で、糖リン酸主鎖の分解ではない。\par

\noindent\textbf{2：不適切} アニーリング温度を上げすぎると非特異的会合が増える。。通常、非特異的な会合は温度が低すぎる場合に増え、高すぎると標的との会合も減り増幅が低下する。製品と配列ごとの条件を確認する。\par

\noindent\textbf{3：不適切} 増幅する配列が長いほど伸長時間を短くする。。同じ伸長速度なら長い領域を合成するにはより長い時間が必要である。\par

\noindent\textbf{4：最適} 高い耐熱性のDNAポリメラーゼを用いる。。各回の加熱変性に耐える酵素を使う。耐熱性と最適反応温度は同じ概念ではない。\par

\noindent\textbf{5：不適切} 生成DNAにプライマーの配列は含まれない。。ポリメラーゼはプライマーの3端から伸長し、プライマー自身は産物の一部に残る。\par

\paragraph{再利用する考え方} 単位・対象・結合の階層を区別する。PCRの3段階と各段階の制御量の成立条件を確認して類似年度の肢順へ対応する。

\paragraph{共通理論} \hyperref[theory:biology:pcr]{第\ref*{theory:biology:pcr}節}, \hyperref[theory:biology:melting]{第\ref*{theory:biology:melting}節}

\par 参考：\cite{neb_pcr}。

\paragraph{関連問題} \hyperref[q:r05-reg-basic-4-6]{R05-REG-BASIC-4-6}

\par 難易度 2（校準初稿）、優先度未評価。

\subsection{令和5年度 問4-6}
\label{q:r05-reg-basic-4-6}

\noindent\textbf{原典：} sources/R05/basic.pdf、PDF物理15頁。

\textbf{公式正答：} 5（single）。

\paragraph{問題の要旨} PCRの3段階、プライマーの産物への組込みと耐熱性を判定する。

\Needspace{4\baselineskip}\paragraph{確認した条件・記述（独自要約）}\mbox{}\par

\noindent\textbf{方向：} 最も適切。\par

\paragraph{問われる核心} PCRの3段階と各段階の制御量

\paragraph{必要知識} PCRの3段階と各段階の制御量。R02の候補順だけでなく耐熱性と熱変性の命題が変わっている。熱によって鎖間会合を解く5が適切。プライマーは産物へ入り、長い標的は伸長時間を要し、低すぎる会合温度は非特異性を増やし得る。

\paragraph{考え方・正答を選ぶ理由（査読待ち初稿）} R02の候補順だけでなく耐熱性と熱変性の命題が変わっている。熱によって鎖間会合を解く5が適切。プライマーは産物へ入り、長い標的は伸長時間を要し、低すぎる会合温度は非特異性を増やし得る。

\Needspace{4\baselineskip}\paragraph{選択肢の検討（査読待ち初稿）}\mbox{}\par

\noindent\textbf{1：不適切} アニーリング温度を上げすぎると非特異的会合が増える。。通常、非特異的な会合は温度が低すぎる場合に増え、高すぎると標的との会合も減り増幅が低下する。製品と配列ごとの条件を確認する。\par

\noindent\textbf{2：不適切} 増幅する配列が長いほど伸長時間を短くする。。同じ伸長速度なら長い領域を合成するにはより長い時間が必要である。\par

\noindent\textbf{3：不適切} 生成DNAにプライマーの配列は含まれない。。ポリメラーゼはプライマーの3端から伸長し、プライマー自身は産物の一部に残る。\par

\noindent\textbf{4：不適切} 低い耐熱性のDNAポリメラーゼがPCRに適する。。反復加熱に耐える必要があり、一般には高い耐熱性を用いる。\par

\noindent\textbf{5：最適} 熱で二重鎖の水素結合を解き一本鎖へ解離させる。。鎖間の塩基対会合を解く熱変性の説明として適切。糖リン酸主鎖の切断や、安定化が水素結合だけとの説明へ拡張しない。\par

\paragraph{再利用する考え方} 単位・対象・結合の階層を区別する。PCRの3段階と各段階の制御量の成立条件を確認して類似年度の肢順へ対応する。

\paragraph{共通理論} \hyperref[theory:biology:pcr]{第\ref*{theory:biology:pcr}節}, \hyperref[theory:biology:melting]{第\ref*{theory:biology:melting}節}

\par 参考：\cite{neb_pcr}。

\paragraph{関連問題} \hyperref[q:r02-reg-basic-4-6]{R02-REG-BASIC-4-6}

\par 難易度 2（校準初稿）、優先度未評価。

\section{タンパク質の配列・折畳み・電荷}

\subsection{平成30年度 問4-6}
\label{q:h30-reg-basic-4-6}

\noindent\textbf{原典：} sources/H30/basic.pdf、PDF物理17頁。

\textbf{公式正答：} 3（single）。

\paragraph{問題の要旨} タンパク質の配列、立体配置、主鎖と側鎖の結合、疎水側鎖の分布を判定する。

\Needspace{4\baselineskip}\paragraph{確認した条件・記述（独自要約）}\mbox{}\par

\noindent\textbf{方向：} 最も適切。\par

\paragraph{問われる核心} タンパク質の配列と折畳みを支える相互作用

\paragraph{必要知識} タンパク質の配列と折畳みを支える相互作用。遺伝子の塩基配列が翻訳時のアミノ酸配列、すなわち一次構造を指定する3が適切。ペプチド結合とジスルフィドは共有結合で、標準の主な配置はLである。疎水側鎖の典型的な内外は水溶性球状タンパク質で論じる。

\paragraph{考え方・正答を選ぶ理由（査読待ち初稿）} 遺伝子の塩基配列が翻訳時のアミノ酸配列、すなわち一次構造を指定する3が適切。ペプチド結合とジスルフィドは共有結合で、標準の主な配置はLである。疎水側鎖の典型的な内外は水溶性球状タンパク質で論じる。

\Needspace{4\baselineskip}\paragraph{選択肢の検討（査読待ち初稿）}\mbox{}\par

\noindent\textbf{1：不適切} 標準アミノ酸はペプチドという非共有結合で連結する。。ペプチド結合はカルボキシ基とアミノ基に由来する共有結合である。\par

\noindent\textbf{2：不適切} タンパク質のアミノ酸はほとんどD体。。リボソームで合成する標準タンパク質は主にL、グリシンは不斉を持たない。\par

\noindent\textbf{3：適切} 一次構造は遺伝子により決まる。。一次構造はアミノ酸配列であり、翻訳される核酸の配列に規定される。\par

\noindent\textbf{4：不適切} 非共有結合として静電相互作用・水素結合・ジスルフィド結合が重要。。ジスルフィド結合は2個のCys側鎖のS-S共有結合であり、非共有という分類が誤り。\par

\noindent\textbf{5：不適切（通常の水溶性球状）} 非極性の側鎖はタンパク質表面に多い。。通常の水溶性球状タンパク質では非極性側鎖は内側へ偏り、水との接触を減らす。膜タンパク質の脂質に接する面等の例外は区別する。\par

\paragraph{再利用する考え方} 単位・対象・結合の階層を区別する。タンパク質の配列と折畳みを支える相互作用の成立条件を確認して類似年度の肢順へ対応する。

\paragraph{共通理論} \hyperref[theory:biology:protein]{第\ref*{theory:biology:protein}節}

\par 参考：\cite{mit_protein}。

\paragraph{関連問題} \hyperref[q:r06-reg-basic-4-6]{R06-REG-BASIC-4-6}

\par 難易度 2（校準初稿）、優先度未評価。

\subsection{令和元年度 問4-6}
\label{q:r01-reg-basic-4-6}

\noindent\textbf{原典：} sources/R01/basic.pdf、PDF物理18頁。

\textbf{公式正答：} 3（single）。

\paragraph{問題の要旨} タンパク質の標準種数・側鎖・主鎖結合と、等電点より高いpHでの電荷を対応させる。

\Needspace{4\baselineskip}\paragraph{確認した条件・記述（独自要約）}\mbox{}\par

\noindent\textbf{空欄：} ア〜エ。水溶液のpHは等電点より高い。\par

\paragraph{問われる核心} タンパク質の側鎖とpH依存電荷

\paragraph{必要知識} タンパク質の側鎖とpH依存電荷。共通の骨格に対してR側鎖が疎水性・酸塩基性等を分ける。アミノ酸間はペプチド結合、標準20種である。等電点pIは正味電荷0のpHであり、pHを上げると解離でプロトンを失う方向となり正味負で、3。pIで各基が全て無電荷になるのではない。

\paragraph{考え方・正答を選ぶ理由（査読待ち初稿）} 共通の骨格に対してR側鎖が疎水性・酸塩基性等を分ける。アミノ酸間はペプチド結合、標準20種である。等電点pIは正味電荷0のpHであり、pHを上げると解離でプロトンを失う方向となり正味負で、3。pIで各基が全て無電荷になるのではない。

\Needspace{4\baselineskip}\paragraph{選択肢の検討（査読待ち初稿）}\mbox{}\par

\noindent\textbf{1：不適切} 15／側鎖／ペプチド／正。標準種数は20で、pHが等電点より高いと正味負に帯電する。\par

\noindent\textbf{2：不適切} 15／アミノ基／エステル／負。種数・各種の物性を分けるR・主鎖結合が違う。負だけが合う。\par

\noindent\textbf{3：最適} 20／側鎖／ペプチド／負。標準20、側鎖の性質、ペプチド、pH>pIで負の全項目が整合。\par

\noindent\textbf{4：不適切} 20／側鎖／エステル／正。種数と側鎖は合うが、連結はペプチド、電荷は負。\par

\noindent\textbf{5：不適切} 20／アミノ基／ペプチド／正。共通のアミノ基を側鎖と取り違え、pHと電荷も逆。\par

\paragraph{再利用する考え方} 単位・対象・結合の階層を区別する。タンパク質の側鎖とpH依存電荷の成立条件を確認して類似年度の肢順へ対応する。

\paragraph{共通理論} \hyperref[theory:biology:amino]{第\ref*{theory:biology:amino}節}, \hyperref[theory:biology:charge]{第\ref*{theory:biology:charge}節}

\par 参考：\cite{mit_protein}。

\par 難易度 2（校準初稿）、優先度未評価。

\subsection{令和6年度 問4-6}
\label{q:r06-reg-basic-4-6}

\noindent\textbf{原典：} sources/R06/basic.pdf、PDF物理16頁。

\textbf{公式正答：} 4（single）。

\paragraph{問題の要旨} タンパク質の配列決定、疎水側鎖の分布、立体配置、共有結合の分類を判定する。

\Needspace{4\baselineskip}\paragraph{確認した条件・記述（独自要約）}\mbox{}\par

\noindent\textbf{方向：} 最も適切。\par

\paragraph{問われる核心} タンパク質の配列と折畳みを支える相互作用

\paragraph{必要知識} タンパク質の配列と折畳みを支える相互作用。核酸の配列が翻訳のアミノ酸配列を指定する4が適切。他は通常の水溶性球状タンパク質の疎水側鎖の位置、LとD、ペプチドとS-Sの共有性を取り違える。H30と類似でも肢順と列挙の範囲を個別に確認する。

\paragraph{考え方・正答を選ぶ理由（査読待ち初稿）} 核酸の配列が翻訳のアミノ酸配列を指定する4が適切。他は通常の水溶性球状タンパク質の疎水側鎖の位置、LとD、ペプチドとS-Sの共有性を取り違える。H30と類似でも肢順と列挙の範囲を個別に確認する。

\Needspace{4\baselineskip}\paragraph{選択肢の検討（査読待ち初稿）}\mbox{}\par

\noindent\textbf{1：不適切（通常の水溶性球状）} 非極性の側鎖はタンパク質表面に多い。。通常の水溶性球状タンパク質では非極性側鎖は内側へ偏り、水との接触を減らす。膜タンパク質の脂質に接する面等の例外は区別する。\par

\noindent\textbf{2：不適切} タンパク質のアミノ酸はほとんどD体。。リボソームで合成する標準タンパク質は主にL、グリシンは不斉を持たない。\par

\noindent\textbf{3：不適切} 標準アミノ酸はペプチドという非共有結合で連結する。。ペプチド結合はカルボキシ基とアミノ基に由来する共有結合である。\par

\noindent\textbf{4：適切} アミノ酸配列は核酸の塩基配列に規定される。。一次構造はアミノ酸配列であり、翻訳される核酸の配列に規定される。\par

\noindent\textbf{5：不適切} 非共有結合として静電・水素・疎水相互作用・ジスルフィド結合が重要。。ジスルフィド結合は2個のCys側鎖のS-S共有結合であり、非共有という分類が誤り。\par

\paragraph{再利用する考え方} 単位・対象・結合の階層を区別する。タンパク質の配列と折畳みを支える相互作用の成立条件を確認して類似年度の肢順へ対応する。

\paragraph{共通理論} \hyperref[theory:biology:protein]{第\ref*{theory:biology:protein}節}

\par 参考：\cite{mit_protein}。

\paragraph{関連問題} \hyperref[q:h30-reg-basic-4-6]{H30-REG-BASIC-4-6}

\par 難易度 2（校準初稿）、優先度未評価。

\section{未確定}

\subsection{平成23年度 問4-1}
\label{q:h23-reg-basic-4-1}

\noindent\textbf{原典：} sources/H23/basic.pdf、PDF物理12頁。

\textbf{公式正答：} 2（single）。

\par\textit{原典本文の校正と独自解説は未作成。分類のみの登録項目。}

\subsection{平成24年度 問4-1}
\label{q:h24-reg-basic-4-1}

\noindent\textbf{原典：} sources/H24/basic.pdf、PDF物理12頁。

\textbf{公式正答：} 1（single）。

\par\textit{原典本文の校正と独自解説は未作成。分類のみの登録項目。}

\subsection{平成24年度 問4-2}
\label{q:h24-reg-basic-4-2}

\noindent\textbf{原典：} sources/H24/basic.pdf、PDF物理12頁。

\textbf{公式正答：} 3（single）。

\par\textit{原典本文の校正と独自解説は未作成。分類のみの登録項目。}

\subsection{平成24年度 問4-4}
\label{q:h24-reg-basic-4-4}

\noindent\textbf{原典：} sources/H24/basic.pdf、PDF物理13頁。

\textbf{公式正答：} 4（single）。

\par\textit{原典本文の校正と独自解説は未作成。分類のみの登録項目。}

\subsection{平成25年度 問4-1}
\label{q:h25-reg-basic-4-1}

\noindent\textbf{原典：} sources/H25/basic.pdf、PDF物理15頁。

\textbf{公式正答：} 2（single）。

\par\textit{原典本文の校正と独自解説は未作成。分類のみの登録項目。}

\subsection{平成25年度 問4-2}
\label{q:h25-reg-basic-4-2}

\noindent\textbf{原典：} sources/H25/basic.pdf、PDF物理15頁。

\textbf{公式正答：} 2（single）。

\par\textit{原典本文の校正と独自解説は未作成。分類のみの登録項目。}

\subsection{平成25年度 問4-4}
\label{q:h25-reg-basic-4-4}

\noindent\textbf{原典：} sources/H25/basic.pdf、PDF物理17頁。

\textbf{公式正答：} 3（single）。

\par\textit{原典本文の校正と独自解説は未作成。分類のみの登録項目。}

\subsection{平成25年度 問4-6}
\label{q:h25-reg-basic-4-6}

\noindent\textbf{原典：} sources/H25/basic.pdf、PDF物理18頁。

\textbf{公式正答：} 3（single）。

\par\textit{原典本文の校正と独自解説は未作成。分類のみの登録項目。}

\subsection{平成26年度 問4-4}
\label{q:h26-reg-basic-4-4}

\noindent\textbf{原典：} sources/H26/basic.pdf、PDF物理14頁。

\textbf{公式正答：} 4（single）。

\par\textit{原典本文の校正と独自解説は未作成。分類のみの登録項目。}

\subsection{平成27年度 問4-3}
\label{q:h27-reg-basic-4-3}

\noindent\textbf{原典：} sources/H27/basic.pdf、PDF物理12頁。

\textbf{公式正答：} 5（single）。

\par\textit{原典本文の校正と独自解説は未作成。分類のみの登録項目。}

\subsection{平成27年度 問4-4}
\label{q:h27-reg-basic-4-4}

\noindent\textbf{原典：} sources/H27/basic.pdf、PDF物理13頁。

\textbf{公式正答：} 1（single）。

\par\textit{原典本文の校正と独自解説は未作成。分類のみの登録項目。}

\subsection{平成28年度 問4-1}
\label{q:h28-reg-basic-4-1}

\noindent\textbf{原典：} sources/H28/basic.pdf、PDF物理17頁。

\textbf{公式正答：} 4（single）。

\par\textit{原典本文の校正と独自解説は未作成。分類のみの登録項目。}

\subsection{平成28年度 問4-3}
\label{q:h28-reg-basic-4-3}

\noindent\textbf{原典：} sources/H28/basic.pdf、PDF物理18頁。

\textbf{公式正答：} 3（single）。

\par\textit{原典本文の校正と独自解説は未作成。分類のみの登録項目。}

\subsection{平成29年度 問4-1}
\label{q:h29-reg-basic-4-1}

\noindent\textbf{原典：} sources/H29/basic.pdf、PDF物理15頁。

\textbf{公式正答：} 4（single）。

\par\textit{原典本文の校正と独自解説は未作成。分類のみの登録項目。}

\subsection{平成29年度 問4-2}
\label{q:h29-reg-basic-4-2}

\noindent\textbf{原典：} sources/H29/basic.pdf、PDF物理15頁。

\textbf{公式正答：} 5（single）。

\par\textit{原典本文の校正と独自解説は未作成。分類のみの登録項目。}

\subsection{平成29年度 問4-3}
\label{q:h29-reg-basic-4-3}

\noindent\textbf{原典：} sources/H29/basic.pdf、PDF物理16頁。

\textbf{公式正答：} 4（single）。

\par\textit{原典本文の校正と独自解説は未作成。分類のみの登録項目。}

\subsection{平成29年度 問4-4}
\label{q:h29-reg-basic-4-4}

\noindent\textbf{原典：} sources/H29/basic.pdf、PDF物理16頁。

\textbf{公式正答：} 1（single）。

\par\textit{原典本文の校正と独自解説は未作成。分類のみの登録項目。}

\subsection{平成30年度 問4-2}
\label{q:h30-reg-basic-4-2}

\noindent\textbf{原典：} sources/H30/basic.pdf、PDF物理15頁。

\textbf{公式正答：} 5（single）。

\par\textit{原典本文の校正と独自解説は未作成。分類のみの登録項目。}

\subsection{平成30年度 問4-3}
\label{q:h30-reg-basic-4-3}

\noindent\textbf{原典：} sources/H30/basic.pdf、PDF物理15頁。

\textbf{公式正答：} 3（single）。

\par\textit{原典本文の校正と独自解説は未作成。分類のみの登録項目。}

\subsection{令和元年度 問4-3}
\label{q:r01-reg-basic-4-3}

\noindent\textbf{原典：} sources/R01/basic.pdf、PDF物理17頁。

\textbf{公式正答：} 4（single）。

\par\textit{原典本文の校正と独自解説は未作成。分類のみの登録項目。}

\subsection{令和元年度 問4-4}
\label{q:r01-reg-basic-4-4}

\noindent\textbf{原典：} sources/R01/basic.pdf、PDF物理17頁。

\textbf{公式正答：} 3（single）。

\par\textit{原典本文の校正と独自解説は未作成。分類のみの登録項目。}

\subsection{令和元年度再試験 問4-1}
\label{q:r01-ret-basic-4-1}

\noindent\textbf{原典：} sources/R01\_retest/basic.pdf、PDF物理14頁。

\textbf{公式正答：} 2（single）。

\par\textit{原典本文の校正と独自解説は未作成。分類のみの登録項目。}

\subsection{令和元年度再試験 問4-2}
\label{q:r01-ret-basic-4-2}

\noindent\textbf{原典：} sources/R01\_retest/basic.pdf、PDF物理14頁。

\textbf{公式正答：} 3（single）。

\par\textit{原典本文の校正と独自解説は未作成。分類のみの登録項目。}

\subsection{令和元年度再試験 問4-3}
\label{q:r01-ret-basic-4-3}

\noindent\textbf{原典：} sources/R01\_retest/basic.pdf、PDF物理14頁。

\textbf{公式正答：} 2（single）。

\par\textit{原典本文の校正と独自解説は未作成。分類のみの登録項目。}

\subsection{令和元年度再試験 問4-4}
\label{q:r01-ret-basic-4-4}

\noindent\textbf{原典：} sources/R01\_retest/basic.pdf、PDF物理15頁。

\textbf{公式正答：} 4（single）。

\par\textit{原典本文の校正と独自解説は未作成。分類のみの登録項目。}

\subsection{令和2年度 問4-2}
\label{q:r02-reg-basic-4-2}

\noindent\textbf{原典：} sources/R02/basic.pdf、PDF物理15頁。

\textbf{公式正答：} 3（single）。

\par\textit{原典本文の校正と独自解説は未作成。分類のみの登録項目。}

\subsection{令和2年度 問4-3}
\label{q:r02-reg-basic-4-3}

\noindent\textbf{原典：} sources/R02/basic.pdf、PDF物理16頁。

\textbf{公式正答：} 4（single）。

\par\textit{原典本文の校正と独自解説は未作成。分類のみの登録項目。}

\subsection{令和2年度 問4-4}
\label{q:r02-reg-basic-4-4}

\noindent\textbf{原典：} sources/R02/basic.pdf、PDF物理17頁。

\textbf{公式正答：} 5（single）。

\par\textit{原典本文の校正と独自解説は未作成。分類のみの登録項目。}

\subsection{令和2年度 問4-5}
\label{q:r02-reg-basic-4-5}

\noindent\textbf{原典：} sources/R02/basic.pdf、PDF物理18頁。

\textbf{公式正答：} 3（single）。

\par\textit{原典本文の校正と独自解説は未作成。分類のみの登録項目。}

\subsection{令和3年度 問4-2}
\label{q:r03-reg-basic-4-2}

\noindent\textbf{原典：} sources/R03/basic.pdf、PDF物理16頁。

\textbf{公式正答：} 4（single）。

\par\textit{原典本文の校正と独自解説は未作成。分類のみの登録項目。}

\subsection{令和3年度 問4-4}
\label{q:r03-reg-basic-4-4}

\noindent\textbf{原典：} sources/R03/basic.pdf、PDF物理18頁。

\textbf{公式正答：} 1（single）。

\par\textit{原典本文の校正と独自解説は未作成。分類のみの登録項目。}

\subsection{令和4年度 問4-1}
\label{q:r04-reg-basic-4-1}

\noindent\textbf{原典：} sources/R04/basic.pdf、PDF物理14頁。

\textbf{公式正答：} 2（single）。

\par\textit{原典本文の校正と独自解説は未作成。分類のみの登録項目。}

\subsection{令和4年度 問4-3}
\label{q:r04-reg-basic-4-3}

\noindent\textbf{原典：} sources/R04/basic.pdf、PDF物理14頁。

\textbf{公式正答：} 1（single）。

\par\textit{原典本文の校正と独自解説は未作成。分類のみの登録項目。}

\subsection{令和5年度 問4-2}
\label{q:r05-reg-basic-4-2}

\noindent\textbf{原典：} sources/R05/basic.pdf、PDF物理13頁。

\textbf{公式正答：} 5（single）。

\par\textit{原典本文の校正と独自解説は未作成。分類のみの登録項目。}

\subsection{令和5年度 問4-3}
\label{q:r05-reg-basic-4-3}

\noindent\textbf{原典：} sources/R05/basic.pdf、PDF物理13頁。

\textbf{公式正答：} 3（single）。

\par\textit{原典本文の校正と独自解説は未作成。分類のみの登録項目。}

\subsection{令和5年度 問4-4}
\label{q:r05-reg-basic-4-4}

\noindent\textbf{原典：} sources/R05/basic.pdf、PDF物理14頁。

\textbf{公式正答：} 1（single）。

\par\textit{原典本文の校正と独自解説は未作成。分類のみの登録項目。}

\subsection{令和6年度 問4-3}
\label{q:r06-reg-basic-4-3}

\noindent\textbf{原典：} sources/R06/basic.pdf、PDF物理15頁。

\textbf{公式正答：} 3（single）。

\par\textit{原典本文の校正と独自解説は未作成。分類のみの登録項目。}

\subsection{令和6年度 問4-4}
\label{q:r06-reg-basic-4-4}

\noindent\textbf{原典：} sources/R06/basic.pdf、PDF物理16頁。

\textbf{公式正答：} 2（single）。

\par\textit{原典本文の校正と独自解説は未作成。分類のみの登録項目。}

\subsection{令和7年度 問4-2}
\label{q:r07-reg-basic-4-2}

\noindent\textbf{原典：} sources/R07/basic.pdf、PDF物理13頁。

\textbf{公式正答：} 3（single）。

\par\textit{原典本文の校正と独自解説は未作成。分類のみの登録項目。}

\subsection{令和7年度 問4-3}
\label{q:r07-reg-basic-4-3}

\noindent\textbf{原典：} sources/R07/basic.pdf、PDF物理13頁。

\textbf{公式正答：} 3（single）。

\par\textit{原典本文の校正と独自解説は未作成。分類のみの登録項目。}
